\documentclass[11pt,letterpaper]{article}
\usepackage[T1]{fontenc}
\usepackage[utf8]{inputenc}
\usepackage{lmodern}
\usepackage[margin=1in]{geometry}
\usepackage{amsmath,amssymb,amsthm,mathtools,mathrsfs}
\usepackage{booktabs,longtable,array,enumitem,xcolor,microtype,xurl}
\usepackage{fancyhdr}
\usepackage[colorlinks=true,linkcolor=blue!42!black,citecolor=blue!42!black,urlcolor=blue!50!black]{hyperref}
\usepackage{bookmark}
\definecolor{ink}{RGB}{27,49,73}
\definecolor{accent}{RGB}{38,81,115}
\pagestyle{fancy}\fancyhf{}
\fancyhead[L]{\small\sffamily Harmonic tails and Gamma identities}
\fancyhead[R]{\small\sffamily ProveIt research continuation}
\fancyfoot[C]{\thepage}
\setlength{\headheight}{14pt}
\setlength{\emergencystretch}{3em}
\setcounter{tocdepth}{2}
\numberwithin{equation}{section}
\theoremstyle{plain}
\newtheorem{theorem}{Theorem}[section]
\newtheorem{lemma}[theorem]{Lemma}
\newtheorem{proposition}[theorem]{Proposition}
\newtheorem{corollary}[theorem]{Corollary}
\newtheorem{conjecture}[theorem]{Conjecture}
\theoremstyle{definition}
\newtheorem{definition}[theorem]{Definition}
\theoremstyle{remark}
\newtheorem{remark}[theorem]{Remark}
\newtheorem{example}[theorem]{Example}
\newcommand{\Q}{\mathbb Q}\newcommand{\R}{\mathbb R}\newcommand{\C}{\mathbb C}
\newcommand{\N}{\mathbb N}\newcommand{\Z}{\mathbb Z}
\newcommand{\dd}{\,\mathrm d}\newcommand{\iu}{\mathrm i}
\newcommand{\NN}{\mathcal N}\newcommand{\HH}{\mathcal H}
\newcommand{\Acal}{\mathcal A}\newcommand{\cA}{\mathcal A}
\newcommand{\cI}{\mathcal I}
\DeclareMathOperator{\Li}{Li}\DeclareMathOperator{\FP}{FP}
\DeclareMathOperator{\CT}{CT}\DeclareMathOperator{\Res}{Res}
\DeclareMathOperator{\rank}{rank}\DeclareMathOperator{\diag}{diag}
\DeclareMathOperator{\Shi}{Shi}\DeclareMathOperator{\wt}{wt}
\DeclareMathOperator{\sgn}{sgn}
\newcommand{\re}{\operatorname{Re}}\newcommand{\im}{\operatorname{Im}}
\newcommand{\fall}[2]{#1^{\underline{#2}}}
\newcommand{\sourcecommit}{446dca9bd65c20e409d1d2561d4ffad408103e2a}
\newcommand{\src}[1]{\nolinkurl{#1}}
\hypersetup{pdftitle={Exact Harmonic-Tail Reconstruction and Variational Gamma Identities},pdfauthor={Research draft prepared with ChatGPT for Vladimir Reshetnikov},pdfsubject={Polylogarithms, harmonic tails, Gamma differential polynomials, zeta and Stieltjes identities}}
\begin{document}
\hypersetup{pageanchor=false}
\begin{titlepage}
\centering
\vspace*{1.0cm}
{\sffamily\color{accent}\Large PROVEIT RESEARCH CONTINUATION\par}
\vspace{1cm}
{\color{ink}\LARGE\bfseries Exact Harmonic-Tail Reconstruction\par}
\vspace{0.3cm}
{\color{ink}\LARGE\bfseries and Variational Gamma Identities\par}
\vspace{0.75cm}
{\large Gap-four determinants, polynomial telescopers,\par
higher resonances, polylogarithms, and Stieltjes coefficients\par}
\vspace{0.9cm}
{\large Prepared for Vladimir Reshetnikov\par}
\vspace{0.2cm}
{Research draft prepared with ChatGPT\par}
\vspace{0.25cm}
{11 October 2026\par}
\vspace{0.8cm}
\begin{minipage}{0.91\textwidth}
\small
\textbf{Repository baseline.} The analysis uses the exact ProveIt revision
\src{446dca9bd65c20e409d1d2561d4ffad408103e2a},
including the manuscript and all six archives then present in
\src{docs/incoming}. The accompanying source manifest records paths,
Git blob identifiers and SHA-256 hashes.

\medskip
\textbf{Scope of the claims.} The report proves the gap-four subfamily of
the optimal moment-prefix conjecture and derives additional exact
identities. The arbitrary unequal-order conjecture and the short
Gaussian $S_6$ and current $S_8$ evaluations remain open. Formal ranks
refer to explicitly chosen coefficient spaces. Numerical diagnostics
are reported separately from proofs. Historical novelty beyond the
inspected sources is not asserted.
\end{minipage}
\vfill
{\small\sffamily Full proofs \enspace / \enspace exact rational certificates
\enspace / \enspace reproducible numerical checks\par}
\end{titlepage}
\pagenumbering{roman}
\hypersetup{pageanchor=true}

\begin{abstract}
This continuation develops four related methods for exact identities.
First, a fourth-order Bernoulli operator proves that the first $q$
centered harmonic-tail moments have full formal rank for every even pair
$p=q-4\ge2$. The remaining boundary determinant is controlled by a
strict ratio-monotonicity argument, not by extrapolating finite ranks.
Second, every finite late-moment combination for arbitrary positive even
orders admits an explicit polynomial-tail telescoping certificate.
Its negative even Laurent coefficients are exactly the zeta coordinates;
cancellation of the positive zeta coordinates gives a rational
evaluation. Third, the Euler variational derivative yields an exact
three-variable equation for the complete $\psi,\psi',\psi''$ sector of
the inverse-Gamma hierarchy. We solve the square and cube sectors in
closed form and give rational coefficient formulas at every order.
Fourth, the next rank-two Gamma-weight resonance is evaluated with all
cutoff and spectral counterterms retained, producing polylogarithmic
generating functions and convergent zeta evaluations. A final Mellin
calculation relates higher spectral derivatives to explicitly subtracted
logarithmic spatial integrals and Stieltjes coefficients. The report
specifies the inherited results, the new proof scopes, integration
recommendations, and further research targets.
\end{abstract}
\tableofcontents
\clearpage
\pagenumbering{arabic}
\section{Research scope, conventions, and dependencies}
\label{sec:scope}

The repository has a substantial body of overlapping reports. A useful
continuation must first fix its comparison point. We use the commit on
the title page, whose last change added six new research bundles.
The canonical manuscript already incorporates many older results;
the six new bundles contain additional theorems that must also be
credited before identifying an open question. The baseline and its
mathematical status are recorded in \cite{Repo,Incoming} and in the
individual reports \cite{BTD,CTD,DHR,DPL,MR,PTI}.

\subsection{What is continued here}

The optimal-prefix question concerns the centered moments of two
ordinary zeta tails. The incoming reports prove the diagonal case,
the adjacent even-order case $p=q-2$, and the families with smaller
order $2,4,$ or $6$; they leave a general unequal-order determinant
unresolved \cite{BTD,DPL,PTI}. Section~\ref{gap:section} proves the
entire next gap $q-p=4$. Thus its first pair outside those earlier
small-order families is $(p,q)=(8,12)$. Section~\ref{tel:section}
then turns late-coordinate relations into explicit telescopers for
all even pairs, whether or not their optimal-prefix determinant is
known to be nonzero.

The inverse-Gamma normal form and the full digamma--trigamma sector
are inherited from \cite{DPL,PTI}. The explicit request in
\cite[research questions 4--5]{PTI} asks for the next mixed sector,
including arbitrary powers of the first two derivative variables.
Section~\ref{var:section} gives an exact recurrence for the entire
sector with largest derivative index at most two. It also proves
closed formulas for all terms containing exactly two or exactly three
copies of $\psi''$. This is a coefficient theorem in a differential
polynomial algebra; its integration consequences retain every endpoint
term.

The rank-two higher-resonance problem is stated explicitly in
\cite[Section 9]{DHR}. Section~\ref{res:section} treats that target,
including the logarithmic square in the hard cutoff and its difference
from the spectral Laurent constant. The logarithmically weighted
finite-part problem in \cite{DPL} is treated in
Section~\ref{log:section}. These are analytic continuations with
specified counterterms, so their normalization is part of each theorem.

\subsection{Ordinary tails, strict sums, and spectral derivatives}

For $s>1$ and $n\ge0$, put
\begin{equation}\label{scope:tails}
 H_n^{(s)}=\sum_{k=1}^n\frac1{k^s},\qquad
 A_s(n)=\zeta(s)-H_n^{(s)}=\zeta(s,n+1).
\end{equation}
For integer $s\ge2$ this gives
\[
 A_s(n)=\frac{(-1)^s}{(s-1)!}\psi^{(s-1)}(n+1).
\]
The even orders used in the tail sections therefore have positive
polygamma normalization. These are standard Hurwitz identities
\cite[Sections 5.15 and 25.11]{DLMF}.

A multiple zeta value uses strictly decreasing indices:
\[
 \zeta(s_1,\ldots,s_d)=
 \sum_{n_1>\cdots>n_d\ge1}
       n_1^{-s_1}\cdots n_d^{-s_d}
\]
where the series converges, with meromorphic continuation when needed.
A prime on a function of a spectral variable denotes differentiation
in that variable. In particular it never denotes differentiation of a
formula known only at integer index values.

We use generalized Stieltjes constants with the convention
\begin{equation}\label{scope:stieltjes}
 \zeta(1+u,a)=\frac1u+
       \sum_{j\ge0}\frac{(-1)^j}{j!}\gamma_j(a)u^j,
 \qquad a>0,
\end{equation}
and $\gamma_j=\gamma_j(1)$. Thus $\gamma_0(a)=-\psi(a)$.

\subsection{Two precise meanings of finite part}

For a function with a finite Laurent-logarithmic singular expansion at
zero, its original-coordinate integral finite part is
\begin{equation}\label{scope:fp}
 \FP\int_0^1f(x)\dd x
 =\CT_{\varepsilon\downarrow0}
             \int_\varepsilon^1f(x)\dd x.
\end{equation}
The symbol $\CT$ selects the constant multiplying
$\varepsilon^0(\log\varepsilon)^0$. The coordinate is $x$ itself.
For a polynomial $H$ in digamma derivatives this convention implies
\begin{equation}\label{scope:boundary}
 \FP\int_0^1\frac{d}{dx}H(\psi(x),\psi'(x),\ldots)\dd x
 =H(\psi(1),\psi'(1),\ldots)
   -[x^0]H(\psi(x),\psi'(x),\ldots).
\end{equation}
This follows directly by taking the constant in the endpoint difference.
A change of coordinate must be evaluated separately.

The centered tail moment is a different, explicitly specified object:
\begin{equation}\label{scope:centered}
 D_{p,q}(s)=\sum_{n\ge0}
       A_p(n)A_q(n)(n+\tfrac12)^{-s},
 \qquad \re s>3-p-q.
\end{equation}
We continue it meromorphically and write
$D_r=D_{p,q}(-2r)$. For positive even $p,q$, the centered asymptotic
expansion has only even powers; the possible poles are at odd
integers, so all these values are regular. Equivalently, subtract
from $(n+\tfrac12)^{2r}A_p(n)A_q(n)$ its polynomial part in
$n+\tfrac12$. The resulting summand is $O(n^{-2})$, and its ordinary
sum equals $D_r$ because
$\zeta(-2j,\tfrac12)=0$ for $j\ge0$. Section~\ref{tel:section}
proves this relation with the finite-cutoff constants made explicit.

\subsection{Status and reproducibility}

The proofs below establish their all-order claims algebraically or
analytically. Exact rational verification files check the matrix
conjugacies, polynomial certificates, normal forms, and displayed
specializations. Independent numerical computations compare different
analytic representations. They are diagnostics, not interval proofs,
unless a check is explicitly described as rational arithmetic.

No assertion of arithmetic independence follows from a formal-rank
statement. In particular, the symbols used as late zeta coordinates
are treated as formal basis labels even though their numerical values
have known multiplicative relations. Likewise, the new recurrence for
Gamma coefficients does not prove that the corresponding mixed moments
form a minimal numerical constant basis.

% Self-contained research fragment; no document preamble.
% All labels use gap:.  Sources: Differential Polygamma/Lerch, Polygamma
% Tail, and Bernoulli--Tornheim
% incoming reports at ProveIt 446dca9bd65c20e409d1d2561d4ffad408103e2a.

\section{The optimal moment prefix for every even-tail gap of four}
\label{gap:section}

The incoming \emph{Differential Polygamma and Lerch Identities} and
\emph{Polygamma Tail Identities} reports prove the optimal consecutive
moment prefix for adjacent even orders, $q-p=2$, and explicitly leave
the next gap $q-p=4$ open. The incoming
\emph{Bernoulli, Tornheim, and Digamma Identities} report proves all
pairs with $p=2,4,6$ and gives a residual determinant of size $p/2$
for arbitrary even pairs. We prove the entire gap-four family. The
first case outside those already proved families is $(p,q)=(8,12)$
\cite{DPL,PTI,BTD}.

The Differential report already reduces a fixed gap to a residual
determinant of size $(q-p)/2$, whose columns satisfy a positive
recurrence. We rederive that reduction with its normalization visible,
then prove the missing strict separation of column ratios. This shows
that the two-column determinant cannot vanish. These are statements
about rational coefficient matrices; no arithmetic independence of
zeta values is asserted. The fixed-gap input is
\cite[\texttt{ar:gap-det}, \texttt{ar:gap-factor}]{DPL}.

\subsection{The inherited moment coordinates and their conventions}

For even integers $p=2a$, $q=2b$, with $1\le a\le b$, write
\[
 A_p(n)=\zeta(p,n+1)=\zeta(p)-H_n^{(p)},\qquad x_n=n+\tfrac12,
\]
and let $D_r$ denote the continued value at $s=-2r$ of
\[
 D_{p,q}(s)=\sum_{n\ge0}A_p(n)A_q(n)x_n^{-s}.
\]
The source results prove regularity at these even arguments and the
inverse Bernoulli transforms
\begin{align}
 D_r&=\sum_{j=0}^r\frac{\binom{2r}{2j}}{2j+1}
 B_{2r-2j}(1/2)C_j,\label{gap:forward}\\
 C_j&=\sum_{r=0}^j\binom{2j+1}{2r}2^{-2(j-r)}D_r.
 \label{gap:inverse}
\end{align}
For $j\ge b$ their exact late coordinates are
\begin{equation}\label{gap:late-coordinates}
 C_j=\sum_{k=0}^{b-1}a_{jk}\zeta(2k),\qquad \zeta(0)=-\tfrac12.
\end{equation}
With $d=2j+1$ and $e=d-p-q+2k+1$, the inherited entries are zero
for $e<0$, and otherwise
\begin{equation}\label{gap:source-entry}
 a_{jk}=B_e\left\{
 \frac{\binom{d-q+1}{p-2k}}{d-q+1}
 +\frac{\binom{d-p+1}{q-2k}}{d-p+1}\right\}.
\end{equation}
Each binomial summand is omitted unless its lower index lies between
$2$ and its upper index. In particular, lower index zero is omitted.
The Bernoulli convention is $B_1=-1/2$.

Set
\begin{equation}\label{gap:late-matrix}
 M_{a,b}=(a_{b+r,k})_{0\le r,k<b}.
\end{equation}
In the formal coordinate order
\[
 C_0,\ldots,C_{b-1},\quad\zeta(0),\zeta(2),\ldots,\zeta(2b-2),
\]
the determinant of the complete first-$q$ moment matrix is
\begin{equation}\label{gap:full-prefix-determinant}
 \frac{\det M_{a,b}}{(4b-1)!!}.
\end{equation}
Indeed the first $b$ rows of the $C$ coordinate matrix form an
identity block, the last $b$ rows have the block $M_{a,b}$, and
the forward transform has diagonal $1/(2r+1)$. These facts and
\eqref{gap:forward}--\eqref{gap:source-entry} are source inputs;
the nonvanishing argument that follows is the new result.

\subsection{A complementary determinant of size half the gap}

Put $m=b-a\ge1$ and $h(t)=\coth(t/2)$. For an even polynomial
$Q(t)=\sum_{n=1}^b q_nt^{2n}$ define
\begin{equation}\label{gap:operator}
 \mathcal L_mQ(t)
 =\bigl(Q^{(2m)}(t)-Q^{(2m)}(0)\bigr)h(t)
                      +(Q(t)h(t))^{(2m)}.
\end{equation}
This is an odd formal Taylor series: the subtraction cancels the
possible pole in the first product, and $Qh$ is an odd Taylor series.
Let $K_{m,b}$ map $(q_1,\ldots,q_b)$ to the coefficients of
$t,t^3,\ldots,t^{2b-1}$ in \eqref{gap:operator}.

\begin{lemma}[Exact operator conjugacy]\label{gap:conjugacy}
If $M_{a,b}^{\mathrm{rev}}$ reverses the columns of $M_{a,b}$, then
\begin{equation}\label{gap:conjugacy-equation}
 K_{m,b}=
 \operatorname{diag}_{0\le r<b}\!\left(\frac2{(2r+1)!}\right)
 M_{a,b}^{\mathrm{rev}}
 \operatorname{diag}_{1\le n\le b}\!\bigl((2n)!\bigr).
\end{equation}
Consequently
\begin{equation}\label{gap:conjugacy-determinant}
 \det K_{m,b}=(-1)^{b(b-1)/2}4^bb!\det M_{a,b}.
\end{equation}
\end{lemma}

\begin{proof}
Write
\[
 h(t)=\sum_{k\ge0}h_kt^{2k-1},\qquad
 h_0=2,\quad h_k=\frac{2B_{2k}}{(2k)!}\quad(k\ge1).
\]
For use in all subsequent formulas define
\begin{equation}\label{gap:AB}
 A_r=\begin{cases}0,&1\le r\le m,\\
 (2r)^{\underline{2m}},&r\ge m+1,
 \end{cases}
 \qquad B_N=(2N-1)^{\underline{2m}}.
\end{equation}
The underline denotes a falling factorial. Row $r$ of $K_{m,b}$
corresponds to $N=r+m+1$. Its entry in column $n$ is
\begin{equation}\label{gap:operator-entry}
 h_{N-n}(A_n+B_N)\quad(n\le N),
\end{equation}
and is zero for $n>N$. The exceptional value $A_m=0$ comes from
the subtraction $Q^{(2m)}(0)$.

To compare this with the source, put $j=b+r$ and $k=b-n$ in
\eqref{gap:source-entry}. The Bernoulli index becomes $2(N-n)$,
while the two upper and lower binomial indices are
\[
 (2r+2,\,2n-2m),\qquad(2r+2m+2,\,2n).
\]
Multiplying by $2(2n)!/(2r+1)!$ gives exactly
\eqref{gap:operator-entry}. The first summand is omitted for
$n\le m$, precisely as required by \eqref{gap:AB}. This proves
\eqref{gap:conjugacy-equation}. The product of its positive diagonal
factors is
\[
 2^b\prod_{n=1}^b\frac{(2n)!}{(2n-1)!}=4^bb!,
\]
and column reversal has sign $(-1)^{b(b-1)/2}$.
\end{proof}

Put
\begin{equation}\label{gap:beta}
 \beta_k=\frac{2|B_{2k}|}{(2k)!}
        =\frac{4\zeta(2k)}{(2\pi)^{2k}}>0\quad(k\ge1),
\end{equation}
and, for $N\ge m+1$, put
\begin{equation}\label{gap:d}
 d_N=2\left\{(2N)^{\underline{2m}}+
                    (2N-1)^{\underline{2m}}\right\}>0.
\end{equation}
Define an array of rational numbers $C_{N,j}^{[m]}$, with
$N\ge1$ and $1\le j\le m$, by
\begin{align}
 C_{r,j}^{[m]}&=\delta_{rj}&& (1\le r\le m),
 \label{gap:initial}\\
 C_{N,j}^{[m]}&=\frac1{d_N}\sum_{r=1}^{N-1}
 \beta_{N-r}(A_r+B_N)C_{r,j}^{[m]}
 &&(N\ge m+1).
 \label{gap:recurrence}
\end{align}
Every entry $C_{N,j}^{[m]}$ with $N\ge m+1$ is strictly positive.
The superscript distinguishes this array from the inherited moment
coordinates $C_j$.

\begin{theorem}[Inherited fixed-gap boundary reduction]
\label{gap:complementary}
For $1\le a<b$ and $m=b-a$,
\begin{equation}\label{gap:complementary-formula}
 \det M_{a,b}
 =\frac{(-1)^{b(b-1)/2}}{4^bb!}
 \left(\prod_{N=m+1}^{b+m}d_N\right)
 \det\bigl(C_{b+i,j}^{[m]}\bigr)_{1\le i,j\le m}.
\end{equation}
Thus the remaining determinant has size $m=(q-p)/2$, independently
of the larger order $q$.
\end{theorem}

\begin{proof}
The sign relation $h_k=(-1)^{k-1}\beta_k$ for $k\ge1$ shows that
\begin{equation}\label{gap:formal-null-solutions}
 Q_j(t)=\sum_{n\ge1}(-1)^{n-j}C_{n,j}^{[m]}t^{2n}
 \qquad(1\le j\le m)
\end{equation}
formally satisfy $\mathcal L_mQ_j=0$. Conversely, the coefficient
at $t^{2N-2m-1}$ determines the next unknown coefficient uniquely:
its nonzero leading factor is $d_N$. Hence these are the unique
formal null solutions with their specified first $m$ coefficients.
This argument requires no convergence of the formal series.

Replace the first $m$ columns of $K_{m,b}$ by the corresponding
linear combinations specified by the truncations of
\eqref{gap:formal-null-solutions} through $n=b$. The column
transformation has determinant one, since those truncations have
initial block equal to the identity. Their first $b-m$ rows vanish.
Their bottom $m$ rows form the block $-TQ_{\mathrm{future}}$, where
$T$ is lower triangular with diagonal $d_{b+1},\ldots,d_{b+m}$,
and
\[
 (Q_{\mathrm{future}})_{ij}
      =(-1)^{b+i-j}C_{b+i,j}^{[m]}.
\]
This follows by writing the omitted terms in each exact formal
coefficient equation; row $b+i$ omits only coefficients
$b+1,\ldots,b+i$.

Move the first $m$ columns to the end. The other diagonal block is
triangular with diagonal $d_{m+1},\ldots,d_b$. Therefore
\[
 \det K_{m,b}
 =(-1)^{m(b-m)+m}\left(\prod_{N=m+1}^{b+m}d_N\right)
                          \det Q_{\mathrm{future}}.
\]
Factoring the row and column signs in $Q_{\mathrm{future}}$
gives $\det Q_{\mathrm{future}}=(-1)^{mb}
\det(C_{b+i,j}^{[m]})$. The combined exponent
$2mb-m(m-1)$ is even. Apply
\eqref{gap:conjugacy-determinant} to obtain
\eqref{gap:complementary-formula}.
\end{proof}

\subsection{Strict separation of the recurrence columns}

The classical Bernoulli Fourier formula, DLMF~24.8.1 at zero,
gives the second equality in \eqref{gap:beta} \cite{DLMF}. In particular, with
$x_\ell=(2\pi\ell)^{-2}$,
\[
 \beta_k=4\sum_{\ell\ge1}x_\ell^k\qquad(k\ge1).
\]
Strict Cauchy--Schwarz implies
\begin{equation}\label{gap:strict-log-convexity}
 \beta_{k+1}^2<\beta_k\beta_{k+2}\qquad(k\ge1),
\end{equation}
since the positive support contains distinct points. Thus
$\rho_k=\beta_{k+1}/\beta_k$ is strictly increasing. This familiar
moment argument is used solely to certify nonzero exact determinants.

\begin{lemma}[Ratio separation for every gap]\label{gap:ratio-lemma}
Let $m\ge2$. For any two initial columns $1\le i<j\le m$, the sequence
\begin{equation}\label{gap:ratio}
 R_N=\frac{C_{N,i}^{[m]}}{C_{N,j}^{[m]}}
       \qquad(N\ge m+1)
\end{equation}
is strictly increasing. Consequently every two-by-two minor of the
positive part of this array, with both its row indices and its
initial-column indices increasing, is strictly negative.
\end{lemma}

\begin{proof}
All relevant entries are positive. Since $A_1=\cdots=A_m=0$,
divide the recurrence by the second column's initial forcing
$B_N\beta_{N-j}/d_N$. Set
\begin{equation}\label{gap:normalized-weights}
 f_N=\frac{\beta_{N-i}}{\beta_{N-j}},\qquad
 w_{N,r}=\frac{\beta_{N-r}}{\beta_{N-j}}
                    \left(1+\frac{A_r}{B_N}\right).
\end{equation}
The exact ratio equation is
\begin{equation}\label{gap:ratio-equation}
 R_N=
 \frac{f_N+\displaystyle\sum_{r=m+1}^{N-1}
                   w_{N,r}C_{r,j}^{[m]}R_r}
      {1+\displaystyle\sum_{r=m+1}^{N-1}
                   w_{N,r}C_{r,j}^{[m]}}.
\end{equation}
By \eqref{gap:strict-log-convexity}, $f_N$ strictly increases: it is
a fixed-length product of consecutive $\rho_k$. For every fixed
$r\ge m+1>j$, the quotient $\beta_{N-r}/\beta_{N-j}$ is the
reciprocal of a fixed-length product of consecutive $\rho_k$,
so it strictly decreases as $N$ increases from $r+1$. Moreover
$A_r>0$ and $B_N$ strictly increases. Thus $w_{N,r}$ strictly
decreases with $N$ for each such fixed $r$.

At $N=m+1$ the sums are empty, giving $R_{m+1}=f_{m+1}$.
For the induction step assume the ratios through $R_{N-1}$ are
strictly increasing. Multiply \eqref{gap:ratio-equation} by its
denominator, subtract $R_{N-1}$ times that denominator, and subtract
the preceding row's corresponding zero identity. The term $r=N-1$
vanishes. The result is
\begin{align}
 &\left(1+\sum_{r=m+1}^{N-1}w_{N,r}C_{r,j}^{[m]}\right)
                 (R_N-R_{N-1})\notag\\
 &\quad=f_N-f_{N-1}
  +\sum_{r=m+1}^{N-2}(w_{N,r}-w_{N-1,r})C_{r,j}^{[m]}
                              (R_r-R_{N-1})>0.
 \label{gap:ratio-increment}
\end{align}
The first term is strictly positive. Each term in the sum is
nonnegative, because both differences there are negative. This
also handles the first induction step, when the sum is empty.
The assertion about minors follows by factoring out the positive
entries in their second column.
\end{proof}

\subsection{The gap-four determinant and the optimal prefix}

Specialize to $m=2$. Write
\[
 u_N=C_{N,1}^{[2]},\qquad v_N=C_{N,2}^{[2]}.
\]
For explicit calculation their common recurrence is
\begin{align}
 &(u_1,u_2)=(1,0),\qquad(v_1,v_2)=(0,1),\notag\\
 &y_N=\frac{\displaystyle\sum_{r=1}^{N-1}
                   \beta_{N-r}(A_r+B_N)y_r}
                {16(N-1)^2(2N-1)(2N-3)}\qquad(N\ge3),
 \label{gap:four-recurrence}\\
 &A_1=A_2=0,\qquad
 A_r=(2r)(2r-1)(2r-2)(2r-3)\quad(r\ge3),\notag\\
 &B_N=(2N-1)(2N-2)(2N-3)(2N-4).\notag
\end{align}
Define the strictly positive rational number
\begin{equation}\label{gap:delta}
 \Delta_b=u_{b+2}v_{b+1}-u_{b+1}v_{b+2}>0.
\end{equation}
Its sign follows immediately from Lemma~\ref{gap:ratio-lemma}.

\begin{theorem}[Every pair $q=p+4$]\label{gap:main}
For every integer $b\ge3$, let $p=2b-4$, $q=2b$. Then
\begin{equation}\label{gap:four-determinant}
 \det M_{b-2,b}
 =(-1)^{b(b-1)/2+1}
 \frac{4^b((b+1)!)^2(2b+3)!!(2b+1)!!}{3b!}\,\Delta_b\ne0.
\end{equation}
The moments $D_0,\ldots,D_{q-1}$ therefore have the full formal
rank $q$. Every later centered moment is an explicit rational
linear combination of these first $q$ moments. The determinant
of their full coordinate matrix is
\eqref{gap:four-determinant} divided by $(4b-1)!!$.
\end{theorem}

\begin{proof}
In \eqref{gap:complementary-formula} the endpoint determinant is
$-\Delta_b$. Here
\[
 d_N=16(N-1)^2(2N-1)(2N-3),
\]
so direct multiplication gives
\[
 \prod_{N=3}^{b+2}d_N
 =\frac{16^b((b+1)!)^2(2b+3)!!(2b+1)!!}{3}.
\]
This proves \eqref{gap:four-determinant}. Its scalar factor is
nonzero and $\Delta_b>0$, establishing nonvanishing. The full
rank assertion follows from \eqref{gap:full-prefix-determinant}.
To reconstruct a later moment, use \eqref{gap:inverse} to obtain
$C_0,\ldots,C_{2b-1}$ from the first $2b$ moments, invert
$M_{b-2,b}$ to recover all the even-zeta coordinates, and then use
\eqref{gap:late-coordinates} and \eqref{gap:forward}. Every step is
finite rational arithmetic.
\end{proof}

\subsection{A compact ordinary identity for the first new pair}

For $(p,q)=(8,12)$ the determinant and first inverse row are
\begin{align}\label{gap:812-data}
 \det M_{4,6}&=\frac{263589503}{309657600},\notag\\
 (M_{4,6}^{-1})_{0,*}
 &=\frac8{790768509}
 (469839,1269227,-11819495,56443363,83672540,12834702).
\end{align}
Thus the displayed integer row applied to $C_6,\ldots,C_{11}$
has value $-790768509/16$.

Define the odd polynomial
\begin{equation}\label{gap:812-K}
 \begin{split}
 K(N)=N^{13}\bigl(&469839+1269227N^2-11819495N^4
                   +56443363N^6\\
                 &+83672540N^8+12834702N^{10}\bigr)
 \end{split}
\end{equation}
and its centered difference
\begin{equation}\label{gap:812-W}
 W(x)=K(x+\tfrac12)-K(x-\tfrac12).
\end{equation}
The resulting ordinary, absolutely convergent identity is
\begin{equation}\label{gap:812-ordinary}
 \begin{split}
 \sum_{n=0}^{\infty}\biggl\{&W(x_n)A_8(n)A_{12}(n)
       -\frac{295198146}{77}x_n^4
       -\frac{10254604347}{154}x_n^2
       -\frac{59115956909}{385}\biggr\}
 =-\frac{790768509}{16}.
 \end{split}
\end{equation}
The summand is $O(x_n^{-2})$. To see this directly, the centered
Hurwitz expansion gives the first three terms of $A_8(n)A_{12}(n)$;
multiplication by the degree-$22$ polynomial $W$ gives precisely
the three displayed nondecaying terms. The inverse Bernoulli
transform identifies its continued sum with the integer combination
in \eqref{gap:812-data}; the removed even powers have continued
sum zero because $\zeta(-2r,1/2)=0$ for $r\ge0$. The accompanying
finite rational telescoping certificate in Section~\ref{tel:812}
supplies a second proof of this individual ordinary identity.

In polygamma notation the tail product is
\[
 A_8(n)A_{12}(n)
   =\frac{\psi^{(7)}(n+1)\psi^{(11)}(n+1)}{7!\,11!}.
\]
The compact difference definition \eqref{gap:812-W} avoids listing
twelve large coefficients while specifying the weight exactly.

\subsection{Verification, source update, and remaining limits}

The supplied exact script independently constructs the source matrix
from its binomial formula, constructs the operator matrix by symbolic
differentiation on monomials, and constructs the two recurrence
solutions. Its $323$ rational checks verify conjugacy for
$3\le b\le12$, the determinant formula for $3\le b\le18$,
formal null-series coefficients through $N=24$, and the finite sign
relations used in the proof. These computations are diagnostics for
finite formulas; the proof above establishes the infinite family.
A second standalone replay checks $45$ complementary determinants
for all $1\le a<b\le10$ and $140$ pairwise ratio minors at gaps
$m=2,3,4,5$. Thus the two tail scripts contain $508$ exact checks
in total; none extrapolates its finite range into a proof.

The optimal-prefix source conjecture can now record $q-p=4$ as
proved, alongside the inherited equal-order, adjacent-order, and
$p=2,4,6$ families. The fixed-gap reduction
\eqref{gap:complementary-formula} is the Differential report's
Proposition~\texttt{ar:gap-det}, rederived here; it retains that
attribution. The new input is Lemma~\ref{gap:ratio-lemma}, which
settles its two-dimensional boundary determinant. No new false
inherited theorem was established in this targeted audit.

For $m\ge3$, pairwise ratio separation does not by itself imply
nonvanishing of the $m$-square endpoint determinant. In fact its
sign need not remain fixed as $b$ changes. The next precise target
is the $m=3$ determinant, corresponding to $q-p=6$, with the
already proved overlapping pairs retained. Further useful targets
are a generating function for the two universal solutions in
\eqref{gap:four-recurrence}, and an exact recurrence or arithmetic
normalization for $\Delta_b$ itself. A general higher-minor sign
claim would require an additional proof; it is not a consequence
of the positivity used here.

% Research note: exact continuation of Differential §09's systematic
% rational-telescoping question, independent of continued moments.
\section{A universal finite telescoper for late harmonic-tail moments}
\label{tel:section}

This section resolves the systematic telescoping question in \cite{DPL}.
The coordinate formulas used in its application are inherited from
\cite{DPL,BTD}; the finite certificate and its limiting calculation
are given here for arbitrary positive even orders.

Let $2\le p\le q$ be even integers, $b=q/2$, and
\[
 A_r(N)=\zeta(r)-H_N^{(r)}=\sum_{k>N}k^{-r}.
\]
Choose finitely many rational coefficients $\lambda_j$, $b\le j\le J$,
and put
\[
 K(X)=\sum_{j=b}^J\lambda_jX^{2j+1},\qquad
 P(X)=K(X+\tfrac12)-K(X-\tfrac12).
\]
In particular $P$ is even. Define Faulhaber polynomials by their values
on nonnegative integers:
\[
 U(N)=\sum_{k=1}^N\frac{K(k)}{k^q},\qquad
 V(N)=\sum_{k=1}^N\frac{K(k)}{k^p}.
\]
The conditions $j\ge b$ ensure that the summands are polynomials in
$k$ with positive odd powers. Thus $U,V\in\mathbb Q[N]$ and
$U(0)=V(0)=0$.

Define the finite even Laurent polynomial
\begin{equation}\label{tel:H}
 H(N)=\frac{K(N)}{N^{p+q}}+
       \frac{U(N-1)}{N^p}+\frac{V(N-1)}{N^q}
      =H_+(N)+\sum_{k=1}^{b-1}\frac{h_k}{N^{2k}},
 \qquad h_0=H_+(0),
\end{equation}
where $H_+$ contains the nonnegative powers.
Let $S(X)$ be the nonnegative-power part of the centered asymptotic
expansion of $P(X)A_p(X-1/2)A_q(X-1/2)$.
This is an explicitly computable even polynomial: one uses only finitely
many terms of
\begin{equation}\label{tel:centered}
 A_r(X-\tfrac12)\sim\frac{X^{1-r}}{r-1}
 +\sum_{\ell\ge1}\frac{B_{2\ell}(1/2)(r)_{2\ell-1}}{(2\ell)!}
      X^{1-r-2\ell}.
\end{equation}
Finally let $R\in\mathbb Q[N]$ be the unique polynomial with
\[
 R(0)=0,\qquad R(N)-R(N-1)=H_+(N)-S(N-\tfrac12).
\]

\begin{theorem}[A finite certificate and its limiting value]
\label{tel:main}
With these definitions,
\begin{equation}\label{tel:F}
 F(N)=K(N)A_p(N)A_q(N)+U(N)A_p(N)+V(N)A_q(N)+R(N)
\end{equation}
satisfies the exact finite identity
\begin{align}\label{tel:finite}
 F(N)-F(N-1)={}&P(N-\tfrac12)A_p(N-1)A_q(N-1)
   -S(N-\tfrac12)-\sum_{k=1}^{b-1}\frac{h_k}{N^{2k}}.
\end{align}
Moreover $F(0)=0$ and $\lim_{N\to\infty}F(N)=-h_0/2$.
Consequently the ordinary absolutely convergent sum has the exact value
\begin{equation}\label{tel:evaluation}
 \boxed{\displaystyle
 \sum_{n=0}^\infty
 \left\{P(n+\tfrac12)A_p(n)A_q(n)-S(n+\tfrac12)\right\}
 =-\frac{h_0}{2}+\sum_{k=1}^{b-1}h_k\zeta(2k).}
\end{equation}
If $h_1=\cdots=h_{b-1}=0$, the certificate consists entirely of
polynomial multiples of the product of the two tails, the two individual
tails, and a polynomial; its value is the rational number $-h_0/2$.
\end{theorem}

\begin{proof}
For a positive odd integer $m$, write
\[
 \sum_{k=1}^N k^m=E_m(N)+\frac12N^m,
 \qquad E_m(-N)=E_m(N),\quad E_m(0)=0.
\]
This follows immediately by expanding
$\{B_{m+1}(N+1)-B_{m+1}\}/(m+1)$: apart from its half-power term,
all surviving powers are positive even powers. Hence there are even
polynomials $U_e,V_e$, both vanishing at zero, such that
\[
 U(N)=U_e(N)+\frac{K(N)}{2N^q},\qquad
 V(N)=V_e(N)+\frac{K(N)}{2N^p}.
\]
Their difference laws imply
$U(N-1)=U_e(N)-K(N)/(2N^q)$, and similarly for $V$.
Therefore
\begin{equation}\label{tel:Hparity}
 H(N)=\frac{U_e(N)}{N^p}+\frac{V_e(N)}{N^q}.
\end{equation}
This proves its evenness. Since $U_e,V_e$ have no constant term, the
most negative possible exponent is $2-q$, which gives the stated range
$1\le k\le b-1$.

Set $a=A_p(N)$ and $b'=A_q(N)$ and use
$A_p(N-1)=a+N^{-p}$ and $A_q(N-1)=b'+N^{-q}$.
Expanding $F(N)-F(N-1)$ and using
\[
 \Delta K=P(N-\tfrac12),\quad
 \Delta U=K(N)N^{-q},\quad \Delta V=K(N)N^{-p}
\]
proves \eqref{tel:finite} as a rational-function identity in
$N,a,b'$. No identity among zeta values enters this step.
The centered expansion \eqref{tel:centered} shows that its first
two terms together are $O(N^{-2})$; the residual even reciprocal
powers are summable. Thus $F(N)$ has a finite limit.

It remains to compute that limit without analytic continuation.
The standard uncentered asymptotic expansion has the parity form
\[
 A_r(N)\sim O_r(N)-\frac12N^{-r},
 \qquad O_r(-N)=-O_r(N),
\]
where $O_r$ is an asymptotic Laurent series containing only odd
negative powers. Combining all three expressions gives the exact
formal asymptotic identity
\[
 KA_pA_q+UA_p+VA_q\sim
 KO_pO_q+U_eO_p+V_eO_q-\frac K4N^{-p-q}
   -\frac{U_e}{2}N^{-p}-\frac{V_e}{2}N^{-q}.
\]
The first four terms have odd powers only. The surviving constant is
\[
 -\frac12\left([N^p]U_e+[N^q]V_e\right)=-\frac{h_0}{2}
\]
by \eqref{tel:Hparity}. The polynomial $R$ has zero constant.
All computations use finitely many asymptotic terms, and the remainder
can be taken $o(1)$ after multiplication by the fixed polynomials.
Since the limit exists, positive powers in this expansion cancel and
the limit equals the displayed constant. Summing \eqref{tel:finite}
and using $F(0)=0$ proves \eqref{tel:evaluation}.
\end{proof}

\paragraph{Exact agreement with the inherited late-block matrix.}
For $d=2j+1$ and $e=d-p-q+2k+1$, the coefficient of $N^{-2k}$ in
the contribution of $N^d$ to $H$ is zero for $e<0$, and otherwise
\[
 a_{jk}=B_e\left\{
 \frac{\binom{d-q+1}{p-2k}}{d-q+1}
 +\frac{\binom{d-p+1}{q-2k}}{d-p+1}\right\}.
\]
Each summand is omitted unless its lower binomial index is between
two and its upper index. Indeed $U_e,V_e$ contain positive even powers
only; the forbidden lower-index-zero term is absent because their
constant terms are zero. Thus $h_k=\sum_j\lambda_j a_{jk}$,
including $k=0$. This is exactly the normalization in
Differential equations \texttt{ar:entries}, \texttt{ar:block}, and
Bernoulli equation \texttt{eq:general-source-entries} \cite{DPL,BTD}.
The theorem supplies the all-order finite certificate requested in
Differential Section 09, ``Systematic rational telescoping for tail
products.'' It proves the certificate for every even pair and every
late-block combination, without presupposing a nonsingular optimal
prefix matrix.

\paragraph{Degree bounds.}
If $\lambda_J\ne0$, then
\[
 \deg K=2J+1,\quad\deg P=2J,\quad
 \deg U=2J+2-q,\quad\deg V=2J+2-p.
\]
Put $\delta=2J+2-p-q$. If $\delta\ge0$, then
$\deg H_+,\deg S\le\delta$ and $\deg R\le\delta+1$.
If $\delta<0$, $H_+=S=R=0$.
For the optimal prefix $J=q-1$ these bounds become
$\deg S\le q-p$ and $\deg R\le q-p+1$.

\subsection{A finite certificate for the (8,12) identity}
\label{tel:812}
For the polynomial $K$ in \eqref{gap:812-K}, set $P=W$ as in
\eqref{gap:812-W}. The exact identity is
\eqref{gap:812-ordinary}, with
$A_8(n)A_{12}(n)=\psi^{(7)}(n+1)\psi^{(11)}(n+1)/(7!\,11!)$.

Here is a compact certificate retaining every rational coefficient.
Put $T=N(N+1)$ and define
\begin{align}\label{tel:812UV}
 U(N)=\frac{T}{8}\bigl(&8556468T^5+32712160T^4-38171739T^3\notag\\
      &+24232272T^2-9577682T+1879356\bigr),\\
 V(N)=\frac{T^5}{8}\bigl(&6417351T^3-3525928T^2-1698454T+1879356\bigr),\notag\\
 R(N)=-\frac{N}{6160}\bigl(&2417202210N^4-5764920315N^3\notag\\
      &+42395905270N^2-129829345800N+260578219219\bigr).
      \label{tel:812R}
\end{align}
The corresponding Laurent polynomial actually has no negative powers:
\[
 H(N)=\frac38\left(4991273N^4+112406360N^2+263589503\right).
\]
Hence $h_0=790768509/8$ and $h_k=0$ for $1\le k\le5$.
Substitution of \eqref{gap:812-K} and \eqref{tel:812UV}--\eqref{tel:812R} into
\eqref{tel:F} gives a finite identity whose difference is exactly
the summand of \eqref{gap:812-ordinary}, and whose limit is the right
side. All its steps are rational-function identities followed by the
parity calculation in the theorem. The accompanying symbolic program
checks the identity in independent variables for $A_8(N),A_{12}(N)$,
and separately checks the limiting constant from finite Laurent
expansions.

\section{A variational equation for all Gamma coefficients of derivative order at most two}
\label{var:section}

This section answers the next-sector question in
\cite[research questions 4--5]{PTI}. The result is independent of the
tail determinant: it is an exact identity in a free differential
polynomial algebra. After substitution of digamma derivatives, it
provides coefficients in identities among Gamma, polygamma, and
height-one zeta integrals.

\subsection{The inherited normal form and the new finite-jet principle}

Let
\[
 \mathscr R=\Q[u_0,u_1,u_2,\ldots],\qquad
 Du_r=u_{r+1},\qquad \wt(u_r)=r+1.
\]
A monomial is \emph{functional} if it is a power of $u_0$, including
$1$, or if its largest positive derivative index occurs at least twice.
Write $\mathscr F$ for their rational linear span.
The classical functional-monomial decomposition
\cite[Proposition 4.10]{Guo2014}, used in \cite{DPL,PTI}, is
\begin{equation}\label{var:normal}
 P=\NN(P)+D\HH(P),\qquad \NN(P)\in\mathscr F,
 \qquad [1]\HH(P)=0.
\end{equation}
It is unique. For completeness, if a monomial $M$ has unique largest
index $r\ge1$, replace $u_r$ by $u_{r-1}$ to form $T$, and let $a$
be the multiplicity of $u_{r-1}$ in $T$. Then
\[
 M=\frac1aDT+\left(M-\frac1aDT\right).
\]
Every remaining monomial has smaller largest derivative index.
This proves existence by termination. If $H$ is nonconstant and has
largest derivative index $r$, the nonzero term
$u_{r+1}\partial H/\partial u_r$ in $DH$ has unique largest index;
it cannot belong to $\mathscr F$. This proves uniqueness. Both maps
preserve factor degree, while $D$ raises weight by one. The quotient
is a vector-space quotient: $\NN$ is linear but is not multiplicative.
For example, $\NN(u_1)=0$ while $\NN(u_1^2)=u_1^2$.

Define the Euler variational derivative
\begin{equation}\label{var:Euler}
 \mathscr E(P)=\sum_{r\ge0}(-D)^r\frac{\partial P}{\partial u_r}.
\end{equation}
The sum is finite for each polynomial. The commutator identity
$\partial_{u_r}D=D\partial_{u_r}+\partial_{u_{r-1}}$, with the
last term zero for $r=0$, proves by cancellation that
$\mathscr E(DH)=0$.

For $d\ge0$ let $\sigma_d$ set all $u_r$, $r>d$, equal to zero,
and let
\[
 \delta_d=\sum_{r=0}^{d-1}u_{r+1}\partial_{u_r}
\]
act on $\Q[u_0,\ldots,u_d]$.

\begin{lemma}[Finite-jet variational principle]\label{var:finitejet}
For every polynomial $P$,
\begin{equation}\label{var:finitejet-eq}
 \sigma_d\mathscr E(P)=
 \sum_{r=0}^d(-\delta_d)^r
      \partial_{u_r}\bigl(\sigma_d\NN(P)\bigr).
\end{equation}
Thus the variational derivative determines a closed equation for a
truncated normal form even though normal reduction itself must precede
truncation.
\end{lemma}
\begin{proof}
Equation \eqref{var:normal} and $\mathscr E D=0$ permit replacing
$P$ by $\NN(P)$ on the left. Any normal monomial removed by
$\sigma_d$ has its highest index greater than $d$ at least twice.
Any partial derivative in \eqref{var:Euler} removes at most one such
factor. Every subsequent application of $D$ preserves at least one
factor of index greater than $d$. Its specialization is therefore zero.
For the retained terms, the ideal generated by $u_{d+1},u_{d+2},\ldots$
is $D$-stable, so specialization of $D^r$ is $\delta_d^r$.
These observations prove the identity term by term.
\end{proof}

\subsection{The inverse-Gamma hierarchy}

Set
\begin{equation}\label{var:hierarchy}
 \Acal_0=1,\qquad \Acal_{m+1}=u_0\Acal_m-D\Acal_m.
\end{equation}
The Bell exponential is
\begin{equation}\label{var:Bell}
 \Acal(t)=\sum_{m\ge0}\Acal_m\frac{t^m}{m!}
 =\exp\left(\sum_{r\ge0}
         \frac{(-1)^ru_rt^{r+1}}{(r+1)!}\right).
\end{equation}
Under $u_r=\psi^{(r)}(x)$, these become
\[
 \Acal_m(x)=(-1)^m\Gamma(x)\frac{d^m}{dx^m}\frac1{\Gamma(x)},
 \qquad \Acal(t,x)=\frac{\Gamma(x)}{\Gamma(x-t)}.
\]
All operations below are formal in $t$; no interchange with an
infinite integral is needed for the coefficient theorem.

Translation $u_0\mapsto u_0+c$ commutes with $D$ and preserves
$\mathscr F$, hence commutes with $\NN$. It follows that
$\NN(\Acal(t))$ has an Appell factor $e^{u_0t}$.
There is also the centered identity
\begin{equation}\label{var:centered}
 \NN(\Acal(t))=
 \NN\exp\left(\sum_{j\ge0}
          \frac{u_{2j}t^{2j+1}}{2^{2j}(2j+1)!}\right).
\end{equation}
Indeed the exponential on the right before normal reduction is
$e^{tD/2}\Acal(t)$, and every coefficient of
$(e^{tD/2}-1)\Acal(t)$ is a total derivative.
Consequently every retained monomial has factor degree congruent to
its weight modulo two.

There are uniquely determined formal series $f_b(y)$ such that
\begin{equation}\label{var:fdefinition}
 \sigma_2\NN(\Acal(t))
 =e^{u_0t} f(y,z),\qquad
 f(y,z)=\sum_{b\ge0}z^bf_b(y),\qquad
 y=\frac{u_1t^2}{2},\quad z=u_2t^3.
\end{equation}
Each $f_b$ is even in $y$. The functional-monomial convention excludes
one copy of $u_2$, so $f_1=0$.

\begin{theorem}[Complete three-variable Gamma-sector equation]
\label{var:master}
Put $L=2y+(z/2)\partial_y$. Then
\begin{equation}\label{var:PDE}
 \boxed{\displaystyle
 f-\frac12L f_y+L^2f_z
 =e^{z/24}\int_0^1
          \cosh(yv)e^{zv^2/8}\dd v.}
\end{equation}
The equation, the normal-form conditions $f_1=0$ and absence of a
linear $y$ term, and the normalization $f_0(0)=1$ uniquely determine
all the functions in \eqref{var:fdefinition}.
More explicitly, let
\begin{equation}\label{var:Rn}
 R_n(y)=\frac1{n!24^n}\int_0^1(1+3v^2)^n\cosh(yv)\dd v.
\end{equation}
Then
\begin{align}
 f_0(y)&=\frac12\left(
       \cosh y+\frac{\sinh y}{y}-y\Shi(y)\right),
       & f_1(y)&=0,\label{var:f0}\\
 4(n+1)y^2f_{n+1}(y)
 &=R_n(y)-(n+1)f_n(y)-(2n-1)y f_n'(y)
           -\frac{n-2}{4}f_{n-1}''(y),
       && n\ge1.\label{var:recurrence}
\end{align}
Here $\Shi(y)=\int_0^y\sinh v\dd v/v$. Every apparent singularity
at $y=0$ in these formulas is removable.
\end{theorem}
\begin{proof}
Differentiating \eqref{var:Bell} with respect to a variable $u_r$
and applying the Euler operator gives
\begin{equation}\label{var:EA}
 \mathscr E\Acal(t)
 =\sum_{r\ge0}\frac{t^{r+1}}{(r+1)!}D^r\Acal(t)
 =\int_0^t e^{sD}\Acal(t)\dd s.
\end{equation}
After setting $u_r=0$ for $r\ge3$, a formal shift replaces $u_0$
by $u_0+s u_1+s^2u_2/2$ and $u_1$ by $u_1+s u_2$.
The exponent of $e^{sD}\Acal(t)$ becomes
\[
 tu_0+(st-t^2/2)u_1+
       (ts^2/2-t^2s/2+t^3/6)u_2.
\]
Substitute $s=t(1+v)/2$. The exponent is
$tu_0+yv+z/24+zv^2/8$. Integrating over $-1\le v\le1$
and using its even part yields the right side of
\eqref{var:PDE}, multiplied by $te^{tu_0}$.

On $e^{tu_0}g(y,z)$ the restricted derivation is
$e^{tu_0}t^{-1}Lg$. Therefore Lemma~\ref{var:finitejet} gives
\[
 \sigma_2\mathscr E\NN\Acal(t)
 =te^{tu_0}\{f-\tfrac12Lf_y+L^2f_z\}.
\]
This proves \eqref{var:PDE}. Expanding
\[
 L^2=4y^2+z+2zy\partial_y+\frac{z^2}{4}\partial_y^2
\]
and comparing $z^n$ gives \eqref{var:recurrence}. At $z^0$ the
condition $f_1=0$ gives
$f_0-yf_0'=\sinh y/y$. Its even solution with constant one is
\eqref{var:f0}; the only homogeneous freedom is a multiple of $y$,
which the normal form excludes.

For $n\ge1$, \eqref{var:recurrence} determines $f_{n+1}$ uniquely
because multiplication by $y^2$ is injective on formal series.
Divisibility of its numerator by $y^2$ follows from the existence of
the normal form already proved in \eqref{var:normal}. This is not an
assumption about arbitrary initial functions. The coefficient theorem
below also proves that each $f_b$ is entire in $y$.
\end{proof}

\subsection{Rational formulas for every coefficient}

For even $k\ge0$, write
\[
 r_b(k)=k![y^k]f_b(y).
\]
Then $r_0(k)=-1/((k-1)(k+1))$ and $r_1(k)=0$.
Taking coefficients in \eqref{var:recurrence} gives the explicit
rational recurrence
\begin{equation}\label{var:rrecurrence}
\begin{split}
 r_{n+1}(k)=\frac1{4(n+1)(k+1)(k+2)}\bigg\{
 &\frac1{n!24^n}\sum_{j=0}^n
              \binom nj\frac{3^j}{k+2j+3}\\
 &-\bigl(n+1+(2n-1)(k+2)\bigr)r_n(k+2)
 -\frac{n-2}{4}r_{n-1}(k+4)\bigg\}.
\end{split}
\end{equation}
This computes a whole coefficient family with rational operations,
without expanding all partitions at its total weight.

\begin{theorem}[Polynomial numerators and fixed odd denominators]
\label{var:rational}
For every $b\ge2$ there is a rational polynomial $p_b(X)$ of degree
exactly $b-2$ such that
\begin{equation}\label{var:rationalform}
 \boxed{\displaystyle
 r_b(k)=\frac{p_b(k)}{\prod_{j=0}^{b}(k+2j+1)},\qquad
 [X^{b-2}]p_b(X)=\frac1{4b!6^{b-1}}.}
\end{equation}
These formulas hold for all even $k\ge0$. In particular every
$f_b(y)$ is entire. If $a,k\ge0$ and $b\ge2$, then
\begin{equation}\label{var:allcoefficients}
 [u_0^au_1^ku_2^b]\NN(\Acal_{a+2k+3b})
 =\begin{cases}
 \displaystyle\frac{(a+2k+3b)!}{a!2^kk!}
       \frac{p_b(k)}{\prod_{j=0}^b(k+2j+1)},& k\text{ even},\\[5pt]
 0,&k\text{ odd}.
 \end{cases}
\end{equation}
A coefficient with one copy of $u_2$ is zero.
\end{theorem}
\begin{proof}
The case $b=2$ follows from \eqref{var:rrecurrence}:
$p_2=1/48$. For a polynomial recurrence, set
\[
 \Delta_n(X)=\prod_{h=0}^n(X+2h+3),\qquad
 S_n(X)=\frac1{n!24^n}\sum_{j=0}^n
       \binom nj3^j\frac{\Delta_n(X)}{X+2j+3}.
\]
Each quotient in this expression is a polynomial. For $n\ge2$,
\eqref{var:rrecurrence} becomes
\begin{equation}\label{var:precurrence}
\begin{split}
 p_{n+1}(X)=\frac1{4(n+1)(X+2)}\bigg\{
 S_n(X)&-\bigl(n+1+(2n-1)(X+2)\bigr)p_n(X+2)\\
 &-\frac{n-2}{4}(X+3)p_{n-1}(X+4)\bigg\},
\end{split}
\end{equation}
where the last term is zero for $n=2$.
The numerator vanishes at $X=-2$: after division by
$\Delta_n(-2)\ne0$, that assertion is exactly the constant-term
identity in \eqref{var:recurrence}, whose left side contains $y^2$.
Thus \eqref{var:precurrence} is a polynomial recurrence, not merely
a rational-function definition.

The degree of $S_n$ is $n$, and its leading coefficient is
$4^n/(n!24^n)=1/(n!6^n)$. The two subtracted terms have degree at
most $n-1$ and $n-2$. Division by $4(n+1)(X+2)$ proves that
$p_{n+1}$ has degree $n-1$ and the stated leading coefficient.
Induction proves \eqref{var:rationalform}.
For fixed $b$, $r_b(k)=O(k^{-3})$ as $k\to\infty$, and the
factorial in the series for $f_b$ proves entire convergence.
Finally substitute $y=u_1t^2/2$, $z=u_2t^3$ in
\eqref{var:fdefinition} and expand $e^{u_0t}$.
The factor-degree parity already proved gives the zero odd coefficients.
\end{proof}

The first polynomial numerators are
\begin{align*}
 p_2(X)&=\frac1{48},&
 p_3(X)&=\frac{2X+5}{1728},\\
 p_4(X)&=\frac{X^2+7X+9}{20736},&
 p_5(X)&=\frac{2X^3+27X^2+100X+90}{1244160},\\
 p_6(X)&=\frac{2X^4+44X^3+316X^2+832X+603}{44789760}.
\end{align*}
The theorem supplies every subsequent numerator by the same finite
recurrence. It does not require that all its displayed possible poles
be actual poles; some factors can cancel at higher $b$.

\subsection{Closed square and cube sectors}

\begin{corollary}[All mixed tetragamma-square and cube coefficients]
\label{var:squarecube}
The two first new entire functions are
\begin{align}
 f_2(y)&=\frac{(y^2+3)\sinh y-3y\cosh y}{48y^5}
       =\frac1{384}\int_0^1(1-v^2)^2\cosh(yv)\dd v,
       \label{var:f2closed}\\
 f_3(y)&=\frac{(2y^4+60y^2+135)\sinh y
                -15y(y^2+9)\cosh y}{1728y^7}.
       \label{var:f3closed}
\end{align}
In particular, for all $a,s\ge0$,
\begin{align}
 [u_0^au_1^{2s}u_2^2]\NN(\Acal_{a+4s+6})
 &=\frac{(a+4s+6)!}
 {48a!2^{2s}(2s)!(2s+1)(2s+3)(2s+5)},
 \label{var:squarecoeff}\\
 [u_0^au_1^{2s}u_2^3]\NN(\Acal_{a+4s+9})
 &=\frac{(4s+5)(a+4s+9)!}
 {1728a!2^{2s}(2s)!(2s+1)(2s+3)(2s+5)(2s+7)}.
 \label{var:cubecoeff}
\end{align}
The removable values are $f_2(0)=1/720$ and $f_3(0)=1/36288$.
\end{corollary}
\begin{proof}
The numerator formulas for $p_2,p_3$ give their power series directly.
For the square sector use
\[
 \int_0^1v^k(1-v^2)^2\dd v
 =\frac8{(k+1)(k+3)(k+5)}.
\]
For the cube sector use
\[
 f_3(y)=\frac{2y\partial_y+5}{82944}
                \int_0^1(1-v^2)^3\cosh(yv)\dd v.
\]
Repeated elementary integration gives \eqref{var:f2closed} and
\eqref{var:f3closed}. Coefficient extraction in
\eqref{var:allcoefficients} proves the final formulas.
\end{proof}

For example, the coefficients of $u_2^3$ in $\NN(\Acal_9)$,
of $u_1^2u_2^2$ in $\NN(\Acal_{10})$, and of $u_2^4$ in
$\NN(\Acal_{12})$ are respectively $10,90,$ and $220$.
At weight $14$, the coefficient of $u_1^4u_2^2$ is $15015$.
These are consequences of the all-order formulas, not the evidence
from which their validity is inferred.

\subsection{What the sector theorem does and does not evaluate}

The specialization $\sigma_2$ is performed after normal reduction.
Substituting $u_r=\psi^{(r)}(x)$ into a full moment identity does not
make the omitted derivatives vanish. For instance
\[
 \NN(\Acal_8)
 =u_0^8-70u_0^4u_1^2+28u_0^2u_2^2-7u_1^4-u_3^2.
\]
The last term remains in the digamma identity. The new theorem solves
the coefficient problem for a complete sector and supplies a
finite-jet variational method for higher sectors; it does not claim
that the sector alone evaluates an arbitrary pure digamma moment.
The next section states the exact transfer with its full boundary
functional.

\section{Exact Gamma moments and their endpoint terms}
\label{mom:section}

The coefficient theorem has an immediate analytic use, but it is
essential to keep the complete normal form and its primitive.
We give the inherited scalar bridge with its normalization, then an
explicit ninth-weight identity containing the first cube term.
The bridge itself is already proved in \cite[Section 4]{PTI}.

\subsection{Height-one Laurent coefficients as ordinary integrals}

Define rational zeta polynomials $c_j$ and height-one Laurent
coefficients $\lambda_d$ by
\begin{equation}\label{mom:clambda}
 \frac1{\Gamma(1+u)}=\sum_{j\ge0}c_ju^j,\qquad
 \lambda_d=[u^1]\zeta(1+u,\{1\}^{d-1}),\qquad d\ge1.
\end{equation}
The coefficient on the right is a Laurent coefficient: for $d>1$
the pole at zero has order $d$. In particular
$c_0=1$, $c_1=\gamma$, and $\lambda_1=-\gamma_1$.
For $v>0$ put
\[
 L(v)=-\log(1-e^{-v}),\qquad
 V_d(v)=L(v)^d-\boldsymbol1_{(0,1)}(v)(-\log v)^d.
\]
Then the normalization can be stated without an unevaluated finite
part:
\begin{equation}\label{mom:lambdaintegral}
 \boxed{\displaystyle
 \lambda_d=c_{d+1}+\frac1{d!}
                    \int_0^\infty\frac{V_d(v)}v\dd v.}
\end{equation}
The integral is ordinary and absolutely convergent. At zero its
integrand is $O((1+|\log v|)^{d-1})$; at infinity it decays
exponentially up to the factor $1/v$.

Here is a direct derivation. With $H_{n-1}(\{1\}^{d-1})$ denoting
the strict unit-index multiple harmonic sum, its generating function
is
\[
 \sum_{n\ge1}H_{n-1}(\{1\}^{d-1})e^{-nv}
 =-\frac1{d!}\frac{d}{dv}L(v)^d.
\]
This follows by differentiating the power series for $L(v)^d$,
or inductively from strict summation. Mellin integration and
integration by parts, initially for $\re u>0$, give
\begin{align}\label{mom:heightmellin}
 \zeta(1+u,\{1\}^{d-1})
 &=\frac1{d!\Gamma(u)}\int_0^\infty v^{u-1}L(v)^d\dd v\notag\\
 &=\frac1{\Gamma(1+u)u^d}
   +\frac{u}{d!\Gamma(1+u)}
          \int_0^\infty v^{u-1}V_d(v)\dd v.
\end{align}
The last integral is holomorphic for $u$ in a neighborhood of zero.
Extracting $[u^1]$ proves \eqref{mom:lambdaintegral} and fixes its sign.

\subsection{The scalar inverse-Gamma evaluation}

Let $\Acal_m$ be the differential polynomials of
\eqref{var:hierarchy}. For $m\ge2$ set
\[
 \mathcal L_m=\FP\int_0^1\Acal_m(x)\dd x,
 \qquad K_{2r}=\frac{(-1)^r\pi^{2r}}{2r+1},\quad K_{2r+1}=0.
\]
The inherited all-order scalar formula is
\begin{equation}\label{mom:scalar}
 \boxed{\displaystyle
 \mathcal L_m=K_m-m!\sum_{d=1}^{m-1}
                (-1)^{m-1-d}c_{m-1-d}\lambda_d.}
\end{equation}
We recall why no coordinate constant is suppressed. Put
\[
 \mathcal C_m(x)=m![t^m]
                  \frac{\Gamma(1-t)\Gamma(x)}{\Gamma(x-t)}.
\]
For $m\ge2$, $\mathcal C_m$ extends continuously to $[0,1]$.
Indeed the residue of the displayed quotient at $x=0$ is exactly
$-t$, with no higher powers of $t$. Euler's beta integral applied to
$F(x,t)=\Gamma(1+t)\Gamma(x)/\Gamma(x+t)$ yields
\begin{align}\label{mom:Phi}
 \int_0^1\{F(x,t)-1+t(\psi(x)+\gamma)\}\dd x
 =t\int_0^\infty
 \frac{(1-e^{-v})^t-1
       -\boldsymbol1_{(0,1)}(v)(v^t-1)}v\dd v.
\end{align}
For clarity, the subtraction follows by integrating in $x$ first
when $\re t>0$. The resulting expression is
\[
 t\int_0^\infty
 \left\{\frac{(1-e^{-v})^t-1}{v}
             +\frac{e^{-v}}{1-e^{-v}}\right\}\dd v-1.
\]
At lower cutoff $\delta$, the second integral is
$-\log(1-e^{-\delta})$, while
$\int_\delta^1(v^t-1)\dd v/v=(1-\delta^t)/t+\log\delta$.
Their limits prove \eqref{mom:Phi}. Its subtracted integrand also
justifies holomorphic continuation and coefficient extraction near
$t=0$.

Replacing $t$ by $-t$ and using \eqref{mom:lambdaintegral} now gives
\[
 \int_0^1\mathcal C_m(x)\dd x
 =-m\int_0^\infty\frac{V_{m-1}(v)}v\dd v
 =m!(c_m-\lambda_{m-1}),\qquad m\ge2.
\]
The constant and linear coefficientwise finite parts of the normalized
quotient are $1$ and $\gamma$, since
$\mathcal C_1=\psi+\gamma$ and $\FP\int_0^1\psi(x)\dd x=0$.
Thus its complete coefficientwise integral is
\[
 \frac1{\Gamma(1+t)}-\sum_{d\ge1}\lambda_dt^{d+1}.
\]
Multiplication by $1/\Gamma(1-t)$ and the reflection identity
$1/(\Gamma(1+t)\Gamma(1-t))=\sin(\pi t)/(\pi t)$ prove
\eqref{mom:scalar}. All uses of an infinite generating series here
are coefficientwise identities justified by the displayed local
subtractions.

\subsection{The complete differential identity and a cube-sector moment}

Write $P_m=\NN(\Acal_m)$ and $H_m=\HH(\Acal_m)$. With
\begin{equation}\label{mom:Delta}
 \Delta H=H(\psi(1),\psi'(1),\ldots)
        -[x^0]H(\psi(x),\psi'(x),\ldots),
\end{equation}
the exact transfer is
\begin{equation}\label{mom:transfer}
 \boxed{\displaystyle
 \FP\int_0^1P_m(\psi,\psi',\ldots)\dd x
 =K_m-m!\sum_{d=1}^{m-1}(-1)^{m-1-d}c_{m-1-d}\lambda_d
       -\Delta H_m.}
\end{equation}
This follows from $\Acal_m=P_m+DH_m$ and
\eqref{scope:boundary}. The endpoint is finite algebra: the standard
local series \cite[Section 5.7]{DLMF} is
\[
 \psi(x)=-x^{-1}-\gamma+
             \sum_{j\ge1}(-1)^{j+1}\zeta(j+1)x^j.
\]
For a monomial $M=\prod u_r^{a_r}$ of weight
$w=\sum(r+1)a_r$, multiply its substitutions by $x^w$ and take
the coefficient of $x^w$. Only the first $w+1$ coefficients of each
normalized series $x^{r+1}\psi^{(r)}(x)$ can contribute. This gives
a terminating algorithm in $\Q[\gamma,\zeta(2),\ldots,\zeta(w)]$.

For example, the full ninth-weight normal polynomial is
\begin{equation}\label{mom:P9}
 P_9=u_0^9-126u_0^5u_1^2+84u_0^3u_2^2
       -63u_0u_1^4-9u_0u_3^2+10u_2^3.
\end{equation}
The coefficient $10$ is the first cube term from
Corollary~\ref{var:squarecube}. For a fully explicit endpoint write
$Z_j=\zeta(j)$ and define
\begin{align}\label{mom:B9}
 B_9={}&-252\gamma^6Z_2+1260\gamma^5Z_3
          +\gamma^4(3150Z_2^2-5670Z_4)\notag\\
 &+\gamma^3(-17640Z_2Z_3+20328Z_5)\notag\\
 &+\gamma^2(-7560Z_2^3+40572Z_2Z_4+16884Z_3^2-57456Z_6)\notag\\
 &+\gamma(28476Z_2^2Z_3-70560Z_2Z_5-54432Z_3Z_4+106416Z_7)\notag\\
 &+1932Z_2^4-22050Z_2^2Z_4-18648Z_2Z_3^2+63036Z_2Z_6\notag\\
 &+47376Z_3Z_5+21942Z_4^2-99624Z_8.
\end{align}
Then $B_9=\Delta H_9$, and the concrete identity is
\begin{align}\label{mom:ninthidentity}
 \FP\int_0^1\bigl\{&\psi^9-126\psi^5(\psi')^2
   +84\psi^3(\psi'')^2-63\psi(\psi')^4
   -9\psi(\psi''')^2+10(\psi'')^3\bigr\}\dd x\notag\\
 &=-9!\sum_{d=1}^8(-1)^{8-d}c_{8-d}\lambda_d-B_9.
\end{align}
All functions inside the integral have argument $x$.
The formula specifies an evaluated combination of mixed moments;
it is not a claim that each term separately reduces to ordinary zeta
values. If desired, every $\lambda_d$ on its right can be replaced
by the convergent integral \eqref{mom:lambdaintegral}.

For comparison, the next full normal polynomial is
\begin{align}\label{mom:P10}
 P_{10}={}&u_0^{10}-210u_0^6u_1^2+210u_0^4u_2^2
       -315u_0^2u_1^4-45u_0^2u_3^2\notag\\
       &+100u_0u_2^3+90u_1^2u_2^2+u_4^2.
\end{align}
The companion exact data file gives $P_m,H_m$ through weight $14$,
both endpoint polynomials at weights $9,10$, and the polynomial
numerators $p_b$ of Theorem~\ref{var:rational} through $b=10$.
The verification program checks the complete differential equality
$\Acal_m=P_m+DH_m$ through weight $24$, not only the retained
three-variable sector.

\subsection{Anchored primitives and higher derivative sectors}

For any real $x,x_0>0$, the same identity gives the ordinary anchored
primitive relation
\begin{equation}\label{mom:anchored}
 \int_{x_0}^x P_m(\psi(t),\psi'(t),\ldots)\dd t
 =\int_{x_0}^x\Acal_m(t)\dd t-H_m(\psi(x),\psi'(x),\ldots)
                           +H_m(\psi(x_0),\psi'(x_0),\ldots).
\end{equation}
It states exactly how the primitive changes under normal reduction;
it does not postulate an elementary antiderivative for $\Acal_m$.
Section~\ref{res:section} gives a separate polylogarithmic family
whose repeated Euler primitives and convergent endpoints are explicit.
For additional differential sectors, Lemma~\ref{var:finitejet} applies
without change with $d=3,4,\ldots$. The resulting equations involve
more variables but retain all variational endpoint information.

% Standalone mathematical notes for integration by the parent agent.
% Source question: incoming DHR, sections/09-research.tex, first subsection.
\section{The next two-particle resonance: all overlap terms and its cutoff constant}
\label{res:sec:resonance}\label{res:section}

The incoming \emph{Determinantal Harmonic Resonance} report
\cite[Section 9, first proposed target]{DHR} proves the
first resonance and explicitly proposes the two-particle next resonance
as its initial higher-resonance target.  We resolve that target here.
The input Gauss summation and the determinant identity are classical;
the contribution relative to that report is a normally convergent
subtraction at the next resonance, including the product overlap, an
explicit hard-cutoff constant, and its exact comparison with the
spectral Laurent constant.  No historical priority claim is intended.

Let $\Re a,\Re c>0$, and put
\[
 b_0=a+c+1,\qquad
 w_n(u)=\frac{\Gamma(n+a)\Gamma(n+c)}
 {\Gamma(n+b_0+u)\Gamma(n+1)},\qquad n\geq0.
\]
The two-particle series
\[
 Z(u)=\sum_{0\leq m<n}(n-m)^2w_m(u)w_n(u)
\]
initially converges for $\Re u>1$.  The next resonance is $u=0$.
Use the abbreviations
\begin{align*}
 t&=ac+a+c, & A&=\frac1{ac}, & h&=\frac{t}{ac}=1+\frac1a+\frac1c,\\
 L(u)&=\frac{\Gamma(a+1)\Gamma(c+1)\Gamma(1+u)}
 {\Gamma(a+1+u)\Gamma(c+1+u)}, &
 A(u)&=\frac{L(u)}{ac},\\
 B&=\psi(1)-\psi(a+1)-\psi(c+1), &
 T&=\psi'(1)-\psi'(a+1)-\psi'(c+1).
\end{align*}
Thus $L(0)=1$, $L'(0)=B$, and $L''(0)=B^2+T$.
For an independent anchor $\Re\alpha>0$, define
\[
 \kappa_\alpha(u)=-t+
       \left(\alpha-a-c-\tfrac12\right)u-\tfrac12u^2.
\]
Powers of $n+\alpha$ use the analytic logarithm on the right
half-plane, real on the positive real axis. For $x>0$, write
$x^{-u}=\exp(-u\log x)$ with the real logarithm.

\begin{theorem}[Normally convergent higher-resonance subtraction]
\label{res:thm:normal}
For $\Re u>-1$, the following three sums converge locally normally,
jointly in $a,c,u,\alpha$ on the stated right-half-plane domains:
\begin{align}
 A(u)&=\sum_{n\geq0}w_n(u),\label{res:eq:A}\\
 B_\alpha(u)&=\sum_{n\geq0}
    \left[nw_n(u)-(n+\alpha)^{-1-u}\right],\label{res:eq:B}\\
 D_\alpha(u)&=\sum_{n\geq0}\left[n^2w_n(u)
    -(n+\alpha)^{-u}
    -\kappa_\alpha(u)(n+\alpha)^{-1-u}\right].\label{res:eq:D}
\end{align}
Every fixed mixed parameter derivative may be taken termwise.  They
have the exact representations
\begin{align}
 B_\alpha(u)&=\frac{L(u)}u-\zeta(1+u,\alpha),\label{res:eq:Bexact}\\
 D_\alpha(u)&=\frac{L(u)(u+t)}{u(u-1)}
    -\zeta(u,\alpha)
    -\kappa_\alpha(u)\zeta(1+u,\alpha).\label{res:eq:Dexact}
\end{align}
Every apparent singularity of these expressions in $\Re u>-1$ is
removable.  The meromorphic continuation of the two-particle sum is
\begin{align}
 Z(u)={}&A(u)\left[\zeta(u,\alpha)
       +\kappa_\alpha(u)\zeta(1+u,\alpha)+D_\alpha(u)\right]
       -\left[\zeta(1+u,\alpha)+B_\alpha(u)\right]^2\notag\\
 ={}&-\frac{L(u)^2(1+u/a)(1+u/c)}{u^2(1-u)}.
 \label{res:eq:Z}
\end{align}
In particular, the square of the subtracted harmonic moment is an
essential overlap term.
\end{theorem}

\begin{proof}
The uniform Gamma-ratio expansion on compact parameter sets
\cite[Section 5.11(iii)]{DLMF} gives
\[
 w_n(u)=n^{-2-u}\left[1+\frac{\beta(u)}n+O(n^{-2})\right],\qquad
 \beta(u)=-t-(a+c+\tfrac12)u-\tfrac12u^2.
\]
When $n$ is replaced by $n+\alpha$ in the leading power of $n^2w_n$,
the next coefficient becomes $\beta(u)+\alpha u=\kappa_\alpha(u)$.
Consequently each summand in \eqref{res:eq:B} and \eqref{res:eq:D}
is $O(n^{-2-\Re u})$, locally uniformly; the same estimate without
subtraction applies to \eqref{res:eq:A}.  On a compact subset of
$\Re u>-1$ these estimates have a summable common majorant.
Cauchy's integral formula on slightly larger compact neighborhoods
then gives the assertion for all fixed mixed derivatives.  In
particular, differentiating at $u=0$ is justified before any
meromorphic expression is evaluated.

For the region in which the relevant ordinary moments converge,
Gauss summation \cite[Equation 15.4.20]{DLMF} yields
\begin{align*}
 \sum_{n\geq0}w_n(u)&=A(u),\\
 \sum_{n\geq0}nw_n(u)&=\frac{L(u)}u,\\
 \sum_{n\geq0}n(n-1)w_n(u)
    &=\frac{(a+1)(c+1)L(u)}{u(u-1)}.
\end{align*}
These formulas follow either by shifting the index in the defining
${}_2F_1$ series, or by differentiating it once and twice in its
argument before setting that argument to one.  Adding the last two
moments gives $L(u)(u+t)/(u(u-1))$.  Subtracting the two Hurwitz model
series proves \eqref{res:eq:Bexact} and \eqref{res:eq:Dexact} first on
nonempty open convergence domains and then by analytic continuation.
The proved holomorphy of the left sides proves removability at both
$u=0$ and $u=1$ wherever they occur in the displayed right sides.

Finite Cauchy--Binet gives
$\sum_{0\leq m<n\leq N}(n-m)^2w_mw_n=M_{0,N}M_{2,N}-M_{1,N}^2$.
Pass to the limit for $\Re u>1$ and insert the preceding moment
formulas.  The numerator simplifies using
$u(u+t)-ac(u-1)=(u+a)(u+c)$, proving \eqref{res:eq:Z}.

There is also a fully summable two-index remainder, rather than only
a formal determinant of continued quantities.  Write $b_n(u)$ and
$d_n(u)$ for the summands in \eqref{res:eq:B} and \eqref{res:eq:D}.
Then
\[
 A(u)D_\alpha(u)-B_\alpha(u)^2
 =\sum_{m,n\geq0}\left[w_m(u)d_n(u)-b_m(u)b_n(u)\right]
\]
is absolutely and locally normally convergent, by the product of
the three $\ell^1$ majorants just proved.  The remaining terms in
\eqref{res:eq:Z} are the explicitly displayed model sums and their
cross terms.  Thus the complete subtraction, including overlaps,
has a normally convergent realization independent of expanding
the Gamma product.
\end{proof}

\begin{theorem}[Hard cutoff and spectral finite part]
\label{res:thm:cutoff}
Let
\[
 S_N=\sum_{0\leq m<n\leq N}(n-m)^2w_m(0)w_n(0),\qquad x=N+1.
\]
As $N\to\infty$, locally uniformly for $\Re a,\Re c>0$,
\begin{equation}
 \boxed{S_N=A x-(\log x+B)^2-h(\log x+B)-2-A
             +O\left(\frac{\log^2(x+1)}x\right).}
 \label{res:eq:cutoff}
\end{equation}
The Laurent expansion of the spectral continuation is
\begin{equation}
 Z(u)=-\frac1{u^2}-\frac{2B+h}u
       -\left[2B^2+T+2hB+h+A\right]+O(u).
 \label{res:eq:Laurent}
\end{equation}
If $C_{\rm cut}$ denotes the constant after subtraction of
$Ax-\log^2x-(2B+h)\log x$, and $C_{\rm sp}$ the Laurent constant,
then
\begin{align}
 C_{\rm cut}&=-B^2-hB-2-A,\label{res:eq:Ccut}\\
 C_{\rm sp}&=2C_{\rm cut}+4+A-h-T.\label{res:eq:compare}
\end{align}
This uses the stated unnormalized spectral deformation $w_n(u)$;
neither constant is being assigned to an unspecified regulator.
\end{theorem}

\begin{proof}
Use the anchor $\alpha=1$ in Theorem~\ref{res:thm:normal}.
The expansion of $L(u)/u$ at zero gives
$B_1(0)=B-\gamma$.  Since $\zeta(0)=-1/2$ and
$\kappa_1'(0)=1/2-a-c$, the constant term in
\eqref{res:eq:Dexact} gives
\[
 D_1(0)=-t(B-\gamma)-ac-1.
\]
Consequently the three finite moments satisfy
\begin{align*}
 M_{0,N}&=A-\frac1x+O(x^{-2}),\\
 M_{1,N}&=\log x+B+O(x^{-1}),\\
 M_{2,N}&=x-t\log x-tB-ac-1+O(x^{-1}).
\end{align*}
The first formula follows by summing the $n^{-2}$ Gamma-ratio
tail; the other two follow from the proved summable residuals and
$H_x=\log x+\gamma+O(x^{-1})$.  Multiplication in the cutoff
determinant gives \eqref{res:eq:cutoff}.  In particular, the product
of the $-1/x$ correction in $M_{0,N}$ with the $x$ divergence in
$M_{2,N}$ contributes $-1$ to the constant.  Omitting it would give
an incorrect finite part.

Finally $L(u)^2=1+2Bu+(2B^2+T)u^2+O(u^3)$ and
\[
 \frac{(1+u/a)(1+u/c)}{1-u}
   =1+hu+(h+A)u^2+O(u^3).
\]
Insert these into \eqref{res:eq:Z}; this proves
\eqref{res:eq:Laurent} and the comparison formula.
\end{proof}

\paragraph{Dependence on the spectral normalization.}
Multiplying each particle's weight by a holomorphic function
$\lambda(u)$ with $\lambda(0)=1$ leaves the hard cutoff at zero
unchanged.  Write $\lambda_j=\lambda^{(j)}(0)$.  Since the total
spectral sum is multiplied by $\lambda(u)^2$, its residue and
constant become, respectively,
\[
 -(2B+h)-2\lambda_1,\qquad
 C_{\rm sp}-2\lambda_1(2B+h)-\lambda_1^2-\lambda_2.
\]
Thus the second derivative of the normalization enters at this
double pole.  These expressions follow by multiplying the two
Taylor--Laurent series; they explain exactly why the first-resonance
rule, which has only one amplitude derivative, cannot be copied
unchanged to the present resonance.

\paragraph{A convergent renormalized sum.}
Let $K_N=S_N-S_{N-1}$, $S_{-1}=0$, and $d=B-\gamma$.
The following is an ordinary absolutely convergent identity:
\begin{equation}
 \sum_{m\geq1}\left[
 K_{m-1}-A+\frac{2H_m+2d+h}{m}-\frac1{m^2}\right]
   =-d^2-hd-2-A.\label{res:eq:mass-sum}
\end{equation}
Indeed the largest-index formula
$K_N=w_N(0)[N^2M_{0,N-1}-2NM_{1,N-1}+M_{2,N-1}]$
and one further Gamma-ratio term give a remainder
$O(\log(m+1)/m^2)$ in the displayed summand.  Its partial sum is
exactly $S_{m-1}-Am+H_m^2+(2d+h)H_m$, and
\eqref{res:eq:cutoff} gives the limit.  The $m=1$ term is included;
the two-particle mass there is zero.

\subsection{Uniform spectral deformation and every cutoff jet}
The previous calculation extends to a uniform formula which retains
the two different escape powers.  In particular, derivatives at the
resonance can be computed without differentiating a nonuniform
asymptotic assertion.

\begin{theorem}[Uniform near-resonance expansion]
\label{res:thm:uniform}
Let $0<\delta<1/3$, $x=N+1$, and
$S_N(u)=\sum_{0\leq m<n\leq N}(n-m)^2w_m(u)w_n(u)$.
Uniformly for $|u|\leq\delta$ and for $a,c$ on compact subsets of
their right half-planes,
\begin{align}
 S_N(u)={}&Z(u)+\frac{A(u)x^{1-u}}{1-u}\notag\\
 &+x^{-u}\left[\frac{2L(u)}{u^2}
              -\frac{A(u)\beta(u)}u-\frac{A(u)}2\right]
       -\frac{x^{-2u}}{u^2(1-u^2)}
       +O\left(x^{-1+2\delta}\log^2(x+1)\right),
 \label{res:eq:uniform}
\end{align}
where $\beta(u)=-t-(a+c+1/2)u-u^2/2$.
The full displayed expression has a removable singularity at $u=0$.
The error is holomorphic in $u$ near this closed disk; every fixed
derivative at $u=0$ tends to zero as $N\to\infty$.
\end{theorem}
\begin{proof}
The uniform Gamma-ratio expansion from
Theorem~\ref{res:thm:normal} and the first two Euler--Maclaurin
terms for the corresponding Hurwitz tails
\cite[Equation 25.11.43]{DLMF} give
\begin{align*}
 M_{0,N}(u)&=A(u)-\frac{x^{-1-u}}{1+u}
                   +O(x^{-2+\delta}),\\
 M_{1,N}(u)&=\frac{L(u)-x^{-u}}u+O(x^{-1+\delta}),\\
 M_{2,N}(u)&=\frac{L(u)(u+t)}{u(u-1)}
       +\frac{x^{1-u}}{1-u}
       -x^{-u}\left[\frac12+\frac{\beta(u)}u\right]
       +O(x^{-1+\delta}).
\end{align*}
For a direct justification of the analytic tails, subtract
$n^{-1-u}$ from $nw_n(u)$ and subtract
$n^{-u}+\beta(u)n^{-1-u}$ from $n^2w_n(u)$ before summing from
$n=x$ to infinity.  The residuals are uniformly summable for
$\Re u>-1$.  Apply the ordinary Hurwitz tail expansion to the
model powers, first in their convergence regions and then by
analytic continuation.  This proves both the error bounds and
their uniform removability at $u=0$.

The apparent quotients involving $u^{-1}$ in the second and third
lines, after their pole cancellations, are
$O(x^\delta\log(x+1))$.  On multiplying the finite moment
determinant, the tail of $M_{0,N}$ times the leading divergence of
$M_{2,N}$ contributes
$-x^{-2u}/((1+u)(1-u))$.  All other terms involving a remainder
or this tail are $O(x^{-1+2\delta}\log^2(x+1))$.
Combine this contribution with the $-x^{-2u}/u^2$ from the square
of $M_{1,N}$; their sum is the last term in
\eqref{res:eq:uniform}.  The other terms simplify to the first
three displayed terms by \eqref{res:eq:Z}.
Finally, use Cauchy's estimate on a circle of radius $\delta$
around zero to transfer the tending-to-zero error to every fixed
derivative.  For clarity, the possible negative Laurent powers in
the full model cancel as follows: their $u^{-2}$ coefficients are
$-1+2-1=0$, and their $u^{-1}$ coefficients are
$-(2B+h)+(2B+h-2\log x)+2\log x=0$.
All other potential poles of the model lie outside $|u|<1$.
Thus the difference between the finite cutoff and the model is
holomorphic throughout $|u|<1$, as required for Cauchy's estimate.
\end{proof}

In particular, for each integer $q\geq0$ there are explicit
polynomials $P_q,Q_q$ of degrees at most $q,q+2$ such that
\begin{equation}
 \left.\partial_u^q S_N(u)\right|_{u=0}
       =xP_q(\log x)+Q_q(\log x)+o(1).
 \label{res:eq:all-cutoff-jets}
\end{equation}
They are obtained by taking $q!$ times the coefficient of $u^q$
in \eqref{res:eq:uniform}; only finitely many Gamma derivatives
are needed at each $q$.  The highest coefficient of the second
polynomial is universal:
\begin{equation}
 [y^{q+2}]Q_q(y)
   =\frac{(-1)^q(2-2^{q+2})}{(q+1)(q+2)}.
 \label{res:eq:universal-log}
\end{equation}
Indeed the only contributions of this degree come from
$2e^{-uy}/u^2-e^{-2uy}/u^2$.  This formula includes the
$-\log^2x$ at $q=0$ and the $+\log^3x$ at $q=1$.
The scale $x^{-2u}$ is essential even when only the finite
part of the undeformed cutoff is wanted.

\subsection{Generalized Stieltjes jets of the subtracted moments}
The normally convergent formula has an auxiliary Hurwitz anchor,
which supplies genuine generalized Stieltjes identities.  This is
a consequence of the subtraction theorem, not a claim that merely
differentiating Gauss summation is new.

Let $Y_q$ be the exponential complete Bell polynomial, with $Y_0=1$,
and set
\[
 \ell_j=\psi^{(j-1)}(1)-\psi^{(j-1)}(a+1)
                    -\psi^{(j-1)}(c+1),\qquad j\geq1.
\]
For $q\geq0$, write
\[
 W_q(n)=w_n(0)Y_q\bigl(-\psi(n+b_0),-\psi'(n+b_0),
                         \ldots,-\psi^{(q-1)}(n+b_0)\bigr).
\]
The interpretation is $W_0(n)=w_n(0)$.  Then
\begin{equation}
 \boxed{\sum_{n\geq0}\left[
 nW_q(n)-\frac{(-\log(n+\alpha))^q}{n+\alpha}\right]
 =\frac{Y_{q+1}(\ell_1,\ldots,\ell_{q+1})}{q+1}
       -(-1)^q\gamma_q(\alpha).}\label{res:eq:Stieltjes}
\end{equation}
All logarithms have the right-half-plane branch.  This is ordinary
absolute convergence, including every fixed mixed derivative in
$a,c,\alpha$.  For example,
\begin{align*}
 \sum_{n\geq0}\left[nw_n(0)-\frac1{n+\alpha}\right]
   &=B+\psi(\alpha),\\
 \sum_{n\geq0}\left[-nw_n(0)\psi(n+b_0)
                 +\frac{\log(n+\alpha)}{n+\alpha}\right]
   &=\tfrac12(B^2+T)+\gamma_1(\alpha).
\end{align*}
These follow from \eqref{res:eq:Bexact} and the defining expansion
$\zeta(1+u,\alpha)=u^{-1}+\sum_{q\geq0}(-1)^q\gamma_q(\alpha)u^q/q!$.
For the second subtracted moment one may use the similarly explicit
coefficient formula
\begin{align}
 \frac{D_\alpha^{(q)}(0)}{q!}
 ={}&-t p_{q+1}-(t+1)\sum_{j=0}^q p_j
     -\frac{\partial_s^q\zeta(0,\alpha)}{q!}\notag\\
 &-\kappa_{q+1}
     -\sum_{j=0}^{\min(2,q)}\kappa_j
                    \frac{(-1)^{q-j}\gamma_{q-j}(\alpha)}{(q-j)!},
 \label{res:eq:Djets}
\end{align}
where $p_j=Y_j(\ell_1,\ldots,\ell_j)/j!$,
$\kappa_0=-t$, $\kappa_1=\alpha-a-c-1/2$,
$\kappa_2=-1/2$, and $\kappa_j=0$ for $j>2$.
The left side is obtained by differentiating the summands of
\eqref{res:eq:D}; thus \eqref{res:eq:Djets} is a finite compiler
for convergent Gamma/polygamma sums in generalized Stieltjes
constants and ordinary spectral Hurwitz-zeta derivatives.

\subsection{A rational specialization and its complete polylogarithm generator}
At $a=c=1$, the weight is $w_n(0)=1/((n+1)(n+2))$.
With $x=N+1$ the elementary finite moments are
\[
 M_{0,N}=\frac{x}{x+1},\quad
 M_{1,N}=H_x+\frac2{x+1}-2,\quad
 M_{2,N}=x-3H_x-\frac4{x+1}+4.
\]
Their determinant simplifies to
\begin{equation}
 \boxed{S_N=x-H_x^2+H_x-1-\frac{H_x-1}{x+1},\qquad
 K_N=1-\frac{H_x}{x}-\frac{H_x-1}{x+1}.}\label{res:eq:rational}
\end{equation}
Thus the cutoff has harmonic degree two, but its largest-index
generator has depth two.  The finite difference removes the constant
coefficient of $H_x^2$; a formal harmonic-degree bound at the cutoff
level need not be sharp for the generator.

For $|z|<1$, summation of \eqref{res:eq:rational} gives
\begin{equation}
 \boxed{\sum_{0\leq m<n}
   \frac{(n-m)^2z^{n+1}}{(m+1)(m+2)(n+1)(n+2)}
 =\frac{z}{1-z}-1-\operatorname{Li}_2(z)
  -\frac{1+z}{2z}\log^2(1-z)-\frac{\log(1-z)}z.}
 \label{res:eq:G}
\end{equation}
The singularities at $z=0$ cancel, and its first term is $z^2/12$.
To check the summation, use
\[
 \sum_{m\geq1}\frac{H_mz^m}{m}
   =\operatorname{Li}_2(z)+\tfrac12\log^2(1-z),\qquad
 \sum_{m\geq1}\frac{H_mz^m}{m+1}
   =\frac{\log^2(1-z)}{2z}.
\]
Both follow by integrating the elementary generating function
$\sum H_mz^m=-\log(1-z)/(1-z)$.
In particular,
\begin{equation}
 \sum_{0\leq m<n}\frac{(n-m)^2\,2^{-n-1}}
 {(m+1)(m+2)(n+1)(n+2)}
 =2\log2-\log^22-\frac{\pi^2}{12}.\label{res:eq:half}
\end{equation}

For any integer $q\geq2$, let
\[
 J_q=\sum_{0\leq m<n}\frac{(n-m)^2}
 {(m+1)(m+2)(n+1)^{q+1}(n+2)}.
\]
It converges absolutely since $K_N\to1$.  We use the strict convention
$\zeta(j,1)=\sum_{m>k\geq1}m^{-j}k^{-1}$.  Put
$E_j=\zeta(j,1)+\zeta(j+1)$ for $j\geq2$,
$R_1=\zeta(2)-1$, and
$R_q=E_q-\zeta(q)-R_{q-1}$ for $q\geq2$.  Then
\begin{equation}
 J_q=\zeta(q)-E_{q+1}-R_q.\label{res:eq:Jall}
\end{equation}
Writing $G(z)$ for \eqref{res:eq:G} and $\Theta=z\,d/dz$, the
corresponding anchored primitives on $|z|<1$ are
\[
 \Theta^{-q}G(z)=\sum_{m\geq1}\frac{K_{m-1}z^m}{m^q}
  =\frac1{(q-1)!}\int_0^z\log^{q-1}(z/t)G(t)\frac{dt}{t}
       \quad(q\geq1),
\]
initially on a positive real segment and then by the analytic power
series.  Their endpoints at one are $J_q$ for $q\geq2$.
Indeed $R_q=\sum_{m\geq1}(H_m-1)/(m^q(m+1))$, and partial
fractions prove its displayed recurrence; the initial value follows
by telescoping.  Euler's classical reduction of $\zeta(j,1)$,
proved in the incoming report \cite[Section 6]{DHR}, gives
a terminating formula using only ordinary zeta values and products
of two such values, with weight at most $q+2$.
Explicitly, for $j\ge2$ that reduction is
\[
 \zeta(j,1)=\frac j2\zeta(j+1)
       -\frac12\sum_{r=1}^{j-2}\zeta(j-r)\zeta(r+1),
\]
with the empty sum interpreted as zero. Hence
\eqref{res:eq:Jall} is a finite ordinary-zeta evaluation at every
integer $q\ge2$.
The first two evaluations are
\begin{align}
 \boxed{J_2=3\zeta(2)-2\zeta(3)-\tfrac54\zeta(4)-1,}
 \label{res:eq:J2}\\
 J_3=1-2\zeta(2)+4\zeta(3)-\tfrac54\zeta(4)
                         -3\zeta(5)+\zeta(2)\zeta(3).
 \label{res:eq:J3}
\end{align}
The primitive $q=1$ diverges logarithmically and is deliberately
excluded from this endpoint statement.

\paragraph{Next questions.}
The next unsolved step after this result is an arbitrary-rank
subtraction at $b-a-c=2r-3$, where a two-by-two divergent moment
block interacts with a nontrivial lower determinant.  A second
target is the full family at $b-a-c=2r-2-M$: separate power and
logarithmic counterterms before comparing their cutoff constants
with spectral Laurent coefficients.  The example above also asks
for a depth invariant which respects finite differencing, rather
than only the degree of a formal cutoff polynomial.

% Primary reference URLs checked by this agent:
% https://dlmf.nist.gov/15.4.E20  (Gauss summation)
% https://dlmf.nist.gov/5.11.E13  (Gamma-ratio asymptotics)
% https://dlmf.nist.gov/25.11.E43 (Hurwitz asymptotics)
% The standalone report should cite these and the pinned incoming DHR report.

% Research continuation of Differential Section 09,
% ``Higher spectral derivatives of the residue calculus.''
\section{Canonical spatial kernels for every spectral coefficient}
\label{log:section}

Let $P\in\mathbb C[u_0,\ldots,u_R]$ be fixed and write
$f_P(x)=P(\psi(x),\psi'(x),\ldots,\psi^{(R)}(x))$.
Its scalar projection is $\phi_P(Y)=P(Y,0,\ldots,0)$.
For a sufficiently large finite $M$, define the principal-part data
\[
 p_j(n)=[v^{-j}]f_P(-n+v),\qquad
 P_0(x)=\sum_{j=1}^M p_j(0)x^{-j},\qquad
 \Delta p_j(n)=p_j(n)-p_j(n-1).
\]
All $p_j(n)$ are explicit polynomials in ordinary harmonic numbers
and zeta values. Define the harmonic Dirichlet functions by
\begin{equation}\label{log:Dirichlet}
 D_j(s)=\sum_{n\ge1}\Delta p_j(n)n^{-s}\qquad(\Re s>0).
\end{equation}
The subscript $j$ denotes pole order, not differentiation order.
We are interested in every Laurent coefficient $[s^q]D_1(s)$,
$q\ge1$. The previously established unweighted finite-part bridge
uses only $q=1$ \cite{DPL,PTI}.

The normalized spatial function and its leading tail are
\begin{equation}\label{log:hK}
 h_P(x)=f_P(x+1)-f_P(x)+P_0(x),\qquad
 K_P(L)=\phi'_P(L)+p_1(0)=\sum_{k\ge0}K_kL^k.
\end{equation}
The $P_0$ term removes the entire singularity at $x=0$.
In particular $h_P(0)=f_P(1)-[x^0]f_P(x)$.

\begin{theorem}[Normalized Mellin identity]
\label{log:Mellin}
The integrals
\begin{equation}\label{log:E}
 E_P(s)=\int_0^1 x^{-s}h_P(x)\,dx
  +\int_1^\infty x^{-s}
       \left\{h_P(x)-\frac{K_P(\log x)}x\right\}\,dx
\end{equation}
define a holomorphic function on $-1<\Re s<1$.
On that strip, as an identity of meromorphic functions,
\begin{equation}\label{log:master}
 \boxed{\displaystyle
 D_1(s)=-\frac{\sin\pi s}{\pi}
       \left(E_P(s)+\sum_{k\ge0}\frac{K_k k!}{s^{k+1}}\right)
 -\sum_{j=2}^M\frac{(s)_{j-1}}{(j-1)!}D_j(j-1+s).}
\end{equation}
The rising factorials are finite polynomials.
Every shifted $D_j$ in the second sum is holomorphic at $s=0$,
with all its derivatives represented by absolutely convergent series.
\end{theorem}

\begin{proof}
The inherited normalized principal-part expansion \cite[Section 2]{DPL}
gives, on
$x>0$, the particularly simple shift identity
\begin{equation}\label{log:shiftres}
 h_P(x)=-\sum_{n\ge1}\sum_{j=1}^M
          \frac{\Delta p_j(n)}{(n+x)^j}.
\end{equation}
It can also be obtained directly by shifting a finite cutoff of the
principal-part expansion: the last boundary term tends to zero since
$p_j(n)=O((1+\log n)^d)$.
The local Laurent expansion proves analyticity at $x=0$.
The differentiable positive-axis polygamma asymptotic expansions give
\[
 f_P(x+1)-f_P(x)=\frac{\phi'_P(\log x)}x
                   +O((1+\log x)^d/x^2),
\]
and hence
$h_P(x)=K_P(\log x)/x+O((1+\log x)^d/x^2)$.
These estimates prove normal convergence of \eqref{log:E}, together
with every fixed $s$-derivative, on its stated strip.

For $0<\Re s<1$, integrate \eqref{log:shiftres} against $x^{-s}$.
The exact beta integral is
\[
 \int_0^\infty\frac{x^{-s}}{(n+x)^j}\,dx
 =n^{1-s-j}\frac{\Gamma(1-s)\Gamma(j-1+s)}{\Gamma(j)}
 =\frac{\pi}{\sin\pi s}\frac{(s)_{j-1}}{(j-1)!}
   n^{-(j-1+s)}.
\]
Fubini is justified by
$\Delta p_j(n)=O((1+\log n)^d/n)$; the slowest case, $j=1$,
is summable for $\Re s>0$ after taking the beta integral.
Also
\[
 \int_1^\infty x^{-s-1}K_P(\log x)\,dx
       =\sum_{k\ge0}K_k k!s^{-k-1}.
\]
Solving for $D_1(s)$ gives \eqref{log:master} on $0<\Re s<1$,
and its right side supplies continuation to the whole strip.
For $j\ge2$, $j-1+s$ has positive real part throughout
$\Re s>-1$, so its series is holomorphic there. This also proves
the assertions about ordinary derivatives.
\end{proof}

For $q\ge1$, define the universal polynomial
\begin{align}\label{log:kernel}
 k_q(L)&=\frac{(-1)^q}{q!}
          \frac{(L+i\pi)^q-(L-i\pi)^q}{2i\pi}\notag\\
 &=(-1)^q\sum_{m=0}^{\lfloor(q-1)/2\rfloor}
       \frac{(-1)^m\pi^{2m}L^{q-2m-1}}
            {(2m+1)!(q-2m-1)!}.
\end{align}
The first four are
\[
 k_1=-1,\qquad k_2=L,\qquad
 k_3=\frac{\pi^2}{6}-\frac{L^2}{2},\qquad
 k_4=\frac{L^3-\pi^2L}{6}.
\]
For any polynomial $w$, put
\[
 \mathcal L_P[w]=\int_0^1w(\log x)h_P(x)\,dx+
 \int_1^\infty w(\log x)
       \left\{h_P(x)-\frac{K_P(\log x)}x\right\}\,dx.
\]
Both integrals are ordinary absolutely convergent integrals.
Set
\begin{align}\label{log:constants}
 A_{P,q}&=-\sum_{\substack{k\ge0\\q+k\ \mathrm{even}}}
        \frac{(-1)^{(q+k)/2}\pi^{q+k}K_k k!}{(q+k+1)!},\notag\\
 T_{P,q}&=\sum_{j=2}^M\sum_{n\ge1}
       \frac{\Delta p_j(n)}{n^{j-1}}
       [s^q]\left\{\frac{(s)_{j-1}}{(j-1)!}e^{-s\log n}\right\}.
\end{align}
The inner coefficient is an explicit polynomial in $\log n$, so
$T_{P,q}$ is an absolutely convergent spectral sum. Equivalently it
is a finite rational combination of the derivatives
$D_j^{(r)}(j-1)$, $0\le r\le q-1$.

\begin{corollary}[Every Laurent coefficient has a normalized spatial kernel]
\label{log:coefficient}
For every fixed differential polynomial $P$ and every integer $q\ge1$,
\begin{equation}\label{log:coefficient-formula}
 \boxed{[s^q]D_1(s)=\mathcal L_P[k_q]+A_{P,q}-T_{P,q}.}
\end{equation}
All terms on the right are fixed by ordinary integrals, convergent
series, and finite endpoint/polynomial calculations.
\end{corollary}
\begin{proof}
Expand $\sin(\pi s)/\pi$ in \eqref{log:master} and differentiate
the normally convergent integrals coefficientwise. The coefficient
of $-\sin(\pi s)x^{-s}/\pi$ is exactly $k_q(\log x)$.
Multiplication with $K_k k!/s^{k+1}$ gives $A_{P,q}$ with the
stated sign and parity. The remaining finite sum gives $T_{P,q}$.
\end{proof}

\subsection{The same scalar projection controls every Stieltjes correction}

Define
\[
 R_P(L)=-\frac{\phi_P(L+i\pi)-\phi_P(L-i\pi)}{2i\pi},
 \qquad R'_P(L)=\sum_{r\ge0}a_rL^r.
\]
The inherited scalar-residue asymptotic \cite[Section 2]{DPL} is
\[
 \Delta p_1(n)=\frac{R'_P(\log n)}n
                 +O((1+\log n)^d/n^2).
\]
With $\zeta(1+s)=s^{-1}+\sum_{m\ge0}(-1)^m\gamma_m s^m/m!$,
the following all-order answer follows.

\begin{corollary}[Complete Stieltjes correction]
\label{log:Stieltjes}
For $q\ge1$,
\begin{equation}\label{log:stieltjes-formula}
 \boxed{\displaystyle
 [s^q]D_1(s)=\frac{(-1)^q}{q!}
 \left\{\sum_{r\ge0}a_r\gamma_{q+r}
 +\sum_{n\ge1}\left(\Delta p_1(n)-\frac{R'_P(\log n)}n\right)
                      (\log n)^q\right\}.}
\end{equation}
The second sum converges absolutely. In particular a mixed polynomial,
for which $\phi_P=0$, requires no Stieltjes correction at any order.
\end{corollary}
\begin{proof}
Subtract the leading polynomial from the Dirichlet series:
\[
 D_1(s)=\sum_r a_r(-1)^r\zeta^{(r)}(1+s)
  +\sum_{n\ge1}\left(\Delta p_1(n)-R'_P(\log n)/n\right)n^{-s}.
\]
The second series is holomorphic for $\Re s>-1$. The coefficient
of $s^q$ in $(-1)^r\zeta^{(r)}(1+s)$ is
$(-1)^q\gamma_{q+r}/q!$. This gives the formula.
\end{proof}

\paragraph{Recovery of the original finite-part normalization.}
For $q=1$, $T_{P,1}=\sum_{j\ge2,n\ge1}
\Delta p_j(n)/((j-1)n^{j-1})=J_P$.
If $C_P=(2\pi)^{-1}\int_{-\pi}^{\pi}\phi_P(i\theta)\,d\theta$,
then direct coefficient extraction gives
$A_{P,1}=\phi_P(0)-C_P$.
Thus the inherited identity
$\operatorname{FP}\int_0^1 f_P=C_P+J_P+[s]D_1(s)$ becomes
\begin{equation}\label{log:fpcheck}
 \operatorname{FP}\int_0^1f_P(x)\,dx=\phi_P(0)-E_P(0).
\end{equation}
This is a normalization check; no endpoint constant has been dropped.
For $q\ge2$, \eqref{log:coefficient-formula} supplies the spatial
normalization requested in Differential Section 09. The choice of the
splitting point $1$ and the full term $K_P(\log x)/x$ is part of the
definition. Altering it without transporting the elementary correction
would change the reported constants.

\subsection{Two explicit logarithmic integral evaluations}

For $P=u_0^2$, the pole differences are
$\Delta p_1(n)=-2/n$, $\Delta p_2(n)=0$, hence
$D_1(s)=-2\zeta(1+s)$ and $T_{P,q}=0$.
Moreover
\[
 h_P(x)=\frac{2(\psi(1+x)+\gamma)}x,\qquad
 K_P(L)=2L+2\gamma.
\]
The case $q=2$ gives the ordinary Stieltjes identity
\begin{equation}\label{log:gamma2}
 \boxed{\displaystyle
 \int_0^1\frac{\log x\,[\psi(1+x)+\gamma]}x\,dx
 +\int_1^\infty\frac{\log x\,[\psi(1+x)-\log x]}x\,dx
 =-\frac{\gamma_2}{2}-\frac{\pi^2\gamma}{6}.}
\end{equation}
The higher $q$ versions follow from the printed universal kernels.
These special-function ingredients have classical antecedents; the
claim here is the all-polynomial normalization theorem, not historical
priority for this isolated scalar integral.

For $P=u_1^2$, the principal parts are
\[
 p_4(n)=1,\quad p_2(n)=2(H_n^{(2)}+\zeta(2)),\quad
 p_1(n)=4(H_n^{(3)}-\zeta(3)),\quad p_3(n)=0.
\]
Thus $D_1(s)=4\zeta(3+s)$, $D_2(s)=2\zeta(2+s)$,
$T_{P,2}=2\zeta'(3)$, and $K_P=-4\zeta(3)$.
Write
\[
 h(x)=-\frac{2}{x^2}
        \bigl(\psi'(1+x)-\zeta(2)+2\zeta(3)x\bigr).
\]
The bracket has a double zero at zero. Equation
\eqref{log:coefficient-formula} gives
\begin{equation}\label{log:trigamma-square}
 \boxed{\displaystyle
 \int_0^1(\log x)h(x)\,dx
 +2\int_1^\infty\frac{\log x}{x^2}
          \bigl(\zeta(2)-\psi'(1+x)\bigr)\,dx
 =2\zeta''(3)+2\zeta'(3)+\frac{2\pi^2}{3}\zeta(3).}
\end{equation}
Both displayed integrals converge absolutely. This example explicitly
tests the origin subtraction, the higher-pole term $T_{P,2}$, and the
$\pi^2$ endpoint correction independently of the scalar Stieltjes case.

\section{Source audit, integration decisions, and further research}
\label{audit:integration-section}

\subsection{The inspected baseline and the scope of the audit}
\label{audit:baseline}

The comparison baseline is the repository revision
\[
 \texttt{446dca9bd65c20e409d1d2561d4ffad408103e2a}.
\]
The two locations named in the research request are the canonical
polylogarithm manuscript and the current incoming directory
\cite{Repo,Incoming}. The six current incoming archives were extracted
and their relevant statements compared with the manuscript's overview
and proof-status records. The targeted mathematical checks concerned
the even-tail coefficient conventions and operator reduction, the
differential-polynomial normal form and known Gamma coefficient sectors,
the normalized residue calculus, and the first determinantal resonance
and its stated continuation problem \cite{DPL,PTI,BTD,DHR}.

The new arguments were checked against those inputs, including the
omitted binomial summand, the signs and index shifts in the boundary
determinant, the Euler-operator projection, and the endpoint constants
in the convergent integral formulas. This is a targeted audit of the
proof chain used here. It is not a certification of every chapter,
historical archive, or earlier numerical calculation in the repository.
\textbf{No new false inherited theorem was established by this audit.}
The source changes below concern proved extensions, attribution, and
precise status updates rather than a newly discovered counterexample.

\subsection{What is new relative to that baseline}
\label{audit:status-matrix}

The following distinctions prevent an already available reduction from
being counted as a new resolution. ``New'' here means established in
this continuation relative to the inspected source material; it does
not assert historical priority over the entire literature.

\begin{center}
\small
\begin{tabular}{@{}>{\raggedright\arraybackslash}p{0.18\linewidth}>{\raggedright\arraybackslash}p{0.34\linewidth}>{\raggedright\arraybackslash}p{0.40\linewidth}@{}}
\hline
Topic & Inherited input & Exact advance in this article\\
\hline
Even-tail rank & Equal and adjacent orders; the $p=2,4,6$ families;
the fixed-gap boundary determinant in DPL,
\texttt{ar:gap-det} \cite{DPL,PTI,BTD}.
& Strict separation of every pair of recurrence-column ratios, and
nonvanishing of the complete $q-p=4$ family. The fixed-gap determinant
is rederived as an input, not claimed as new.\\[4pt]
Finite telescopers & The Bernoulli coordinate formula and explicit
low-order telescopers, including $(4,6)$ \cite{DPL}.
& A polynomial certificate for every late-block combination and every
even pair, with a uniform construction, degree bounds, and limiting
constant. The $(8,12)$ rational extraction is an explicit application.\\[4pt]
Gamma coefficients & The canonical quotient, Appell transport, the
complete $u_0,u_1$ sector, and every pure polygamma-square coefficient
\cite{DPL,PTI,BTD}.
& The complete projection onto $u_0,u_1,u_2$, including all powers
$u_2^b$, $b\ge2$, and the rational coefficient form with numerator
degree $b-2$. These are formal coefficient identities.\\[4pt]
Next resonance & Classical Gauss summation and the determinant
identity; the first-resonance theory and proposed next two-particle
target \cite{DHR}.
& The next double pole, including overlap subtraction, the hard-cutoff
constant, its spectral comparison, uniform cutoff jets, and anchored
Hurwitz--Stieltjes identities.\\[4pt]
Higher spectral jets & The normalized principal-part expansion,
scalar projection, and first spectral-coefficient bridge \cite{DPL}.
& Canonical spatial kernels for every $[s^q]D_1(s)$, $q\ge1$, with
all origin, infinity, and higher-pole terms. The same scalar projection
controls every Stieltjes correction.\\[4pt]
Family--2 example & An exact incoming proof of the manuscript's
\texttt{double:eq:fam2-example}, with a simplified value \cite{PTI}.
& An editorial promotion supported by that inherited proof. No new
proof of this example is claimed here.\\
\hline
\end{tabular}
\end{center}

\subsection{Concrete integration decisions and limits}
\label{audit:decisions}

The optimal-prefix conjecture should now record the union of the
proved equal-order, adjacent-order, $p=2,4,6$, and gap-four families.
Its general unequal-order statement remains open. In particular,
pairwise ratio separation does not settle a boundary determinant of
size three or larger. There is a simple algebraic reason: the positive
matrix $B_{ij}=i+j$, $1\le i,j\le3$, has every increasing
two-by-two minor strictly negative,
\[
 \det B_{\{i_1,i_2\},\{j_1,j_2\}}
       =-(i_2-i_1)(j_2-j_1)<0,
\]
but its three-by-three determinant is zero. Thus the new ratio theorem
supplies exactly the sign information needed for a two-column
boundary block; an additional argument is required for $m\ge3$.

The rational telescoper theorem changes the status of the requested
all-order certificate construction. It does not by itself prove that
an arbitrary prescribed inverse row exists. When a coefficient relation
has been supplied, the theorem constructs and verifies its finite
certificate; the separate rank theorem proves existence of every
reconstruction row in the stated families.

For the Gamma sector, retain the order of operations: first take the
canonical normal form and then apply the projection $\sigma_2$ which
sets $u_r=0$ for $r\ge3$. The Euler-operator proof justifies extracting
this projected sector; setting higher derivatives to zero before
normal reduction is a different operation. The all-$b$ recurrence and
rational numerator theorem extend the earlier sector results. They
neither evaluate every associated finite-part moment nor establish
independence of the resulting numerical constants \cite{DPL,PTI}.

All normalizations should travel with their formulas. The tail
identities use the centered coordinate $n+1/2$ and explicit polynomial
subtraction. Differential-polynomial finite parts retain the original
spatial coordinate and the endpoint functional of total derivatives.
The logarithmic kernels use the displayed subtraction at infinity and
the splitting point $1$. The resonant hard cutoff uses $x=N+1$, while
its spectral constant belongs to the explicitly stated Gamma
deformation. The product of a vanishing tail with a divergent moment
contributes to that constant and must be retained. These distinctions
are requirements of the proved identities, not interchangeable
definitions of one unspecified regularization.

The canonical Family--2 display
\texttt{chapters/04-depth.tex}, \texttt{double:eq:fam2-example},
can be marked proved by citing the existing incoming component
evaluations \texttt{bip:eq:12}--\texttt{bip:eq:21} and simplification
\texttt{bip:eq:family2-short} \cite{PTI}. This promotion was already
justified before the present work. The Gaussian $S_6$ and the
\emph{current} $S_8$ candidate remain conjectural. The older rejected
$S_8$ vector is a different statement and must not replace the current
candidate in either a proof attempt or a status table \cite{Repo,PTI}.
None of the new rank, Gamma-sector, or resonance results proves those
period identities.

Exact finite algebra and numerical diagnostics serve different
purposes. The rational replays test coefficient transformations and
finite certificates. The separately subtracted quadratures and
complex-parameter comparisons test analytic formulas and endpoint
constants. The infinite-family results rest on the proofs, rather than
on the extent of a computed table; the floating-point diagnostics are
not interval certificates. No proof-assistant formalization is claimed.

\subsection{Twelve concrete questions for the next continuation}
\label{audit:future}

The following targets emphasize exact identities, structural reductions,
and normalized primitives. Each calls for a proof or a specified finite
certificate; a failed numerical relation search would not settle an
independence or nonmembership claim.

\begin{enumerate}
\item \textbf{Settle the next boundary determinant.}
For $m=3$, prove or disprove nonvanishing of
$\det(C_{b+i,j}^{[3]})_{1\le i,j\le3}$ for every $b>3$.
This is the next gap $q-p=6$. The pairwise signs are now known, but
they do not control the three-column determinant. A discrete Wronskian
identity, an exact recurrence for the determinant, or a nonvanishing
arithmetic normalization would supply the missing step.

\item \textbf{Explain the gap-four determinant sequence arithmetically.}
Find a generating function, finite recurrence, or exact denominator
formula for $\Delta_b=u_{b+2}v_{b+1}-u_{b+1}v_{b+2}$.
The recurrence already proves positivity; the new target is a closed
identity or an integer normalization which explains the numerators
of the optimal-prefix determinants and supports exact reconstruction.

\item \textbf{Describe all reconstruction rows without generic inversion.}
For the proved gap-four family, express each inverse row through
finite binomial or Bernoulli sums and the two fundamental recurrence
solutions. Combine it with the universal telescoper to obtain
certificates for prescribed even-zeta coordinates. Determine minimal
weight-polynomial degrees within the specified formal coefficient
system, keeping this separate from arithmetic independence.

\item \textbf{Extend polynomial telescoping to three or more tails.}
Construct a finite difference certificate involving products of every
required proper subset of the tails. Identify precisely which parity
conditions leave only even reciprocal powers and which additional
Euler-sum coordinates survive. The target is an all-order certificate
and its endpoint constant, rather than an unproved extrapolation of
the two-tail cancellation.

\item \textbf{Determine the complete next Gamma projection.}
Apply the Euler-operator method to $\sigma_3$, retaining
$u_0,u_1,u_2,u_3$. Derive a closed coefficient equation with the
canonical support restrictions, then solve the first coupled families
containing $u_3^2$ and arbitrary powers of $u_1,u_2$.
Prove that discarded higher canonical monomials cannot contribute
before extracting this sector.

\item \textbf{Find a closed formula for the all-$b$ sector numerators.}
The rational form, denominator, degree, and leading coefficient of
the polynomials $p_b(k)$ are now established. Seek a finite formula
uniform in both $b$ and $k$, a bivariate generating function, or a
hypergeometric representation. Such a formula should explain exact
vanishing coefficients and coefficient patterns, rather than merely
reproduce the already proved recurrence.

\item \textbf{Evaluate the numerical image of the new Gamma sector.}
Map the moments containing $(\psi')^{2k}(\psi'')^b$ and their
Gamma--Appell transports into convergent free-order multiple-zeta
derivatives. Determine the relations among these images within an
explicitly stated functional-relation family. This would connect the
new formal coefficient theorem with the remaining basis-reduction questions for
higher mixed moments in the incoming reports \cite{DPL,PTI}.

\item \textbf{Introduce rational shifts into the logarithmic kernels.}
Replace the unshifted pole lattice by a fixed Hurwitz lattice and
derive the exact analogue of the scalar correction for
$\gamma_q(a)$. At rational $a$, seek finite formulas in a declared
multiple-Hurwitz or cyclotomic basis. Retain the origin subtraction,
the choice of logarithm, and the transport of the normalization
constant under changes of base point.

\item \textbf{Construct parameter primitives for the higher jets.}
Starting with the normally convergent kernels, integrate in a shift
or Gamma parameter and anchor the primitive at a specified point.
Find recurrences which isolate one prescribed higher Stieltjes
coefficient or one marked spectral derivative. A complete answer
should prove the endpoint limit and identify any Barnes-type
functions needed for the resulting primitive family.

\item \textbf{Treat higher-rank and deeper determinantal resonances.}
At the next general-rank resonance $b-a-c=2r-3$, construct the
subtraction for the divergent two-by-two moment block together with
the finite lower block. Then study the deeper loci
$b-a-c=2r-2-M$. Determine all overlapping escape terms before
comparing hard-cutoff constants with spectral Laurent coefficients,
and prove uniformity sufficient for every fixed parameter jet.

\item \textbf{Find a depth description compatible with finite differences.}
For rational Gamma weights at other positive integer parameters,
derive complete polylogarithm generators of the largest-index
increments and their anchored Euler primitives. Determine which
cutoff harmonic products cancel under finite differencing and when
the endpoint sums reduce to ordinary zeta products. The two-particle
specialization here supplies a test case for a depth invariant that
respects this operation.

\item \textbf{Return to the frozen Gaussian period targets.}
For the exact $S_6$ vector and current $S_8$ candidate, seek a
parameter-dependent polylogarithm or iterated-integral identity whose
proved boundary value is the required relation. Preserve the branch
data and all lower-weight constants. A projection outside a selected
formal relation space is not a disproof of the numerical period
identity, and a small residual is not its proof \cite{Repo}.
\end{enumerate}

% Bibliographic facts checked against the six incoming TeX title pages.
% Incoming entries deliberately have no invented personal/corporate author.
% The URLs identify the archives at the current audit snapshot, not the
% older source baselines printed inside the individual reports.
\begin{thebibliography}{10}
\raggedright

\bibitem{Repo}
ProveIt Contributors,
\emph{Polylogarithms and their Arithmetic Bridges:
Special Values, Zeta Jets, Lattices and Experimental Discovery}.
Consolidated manuscript, October 2026.
Repository snapshot
\texttt{446dca9bd65c20e409d1d2561d4ffad408103e2a}.
\href{https://github.com/VladimirReshetnikov/ProveIt/tree/446dca9bd65c20e409d1d2561d4ffad408103e2a/Analysis/Polylogarithms/docs/manuscript}
{\nolinkurl{Analysis/Polylogarithms/docs/manuscript}}.

\bibitem{Incoming}
\emph{Incoming reports}.
ProveIt repository inventory and intake instructions, at snapshot
\texttt{446dca9bd65c20e409d1d2561d4ffad408103e2a}.
\href{https://github.com/VladimirReshetnikov/ProveIt/tree/446dca9bd65c20e409d1d2561d4ffad408103e2a/docs/incoming}
{\nolinkurl{docs/incoming}};
\href{https://raw.githubusercontent.com/VladimirReshetnikov/ProveIt/446dca9bd65c20e409d1d2561d4ffad408103e2a/docs/incoming/README.md}
{pinned intake README}.
The six archives below are repository working reports; their presence
in the incoming directory does not confer canonical or formal status.

\bibitem{BTD}
\emph{Bernoulli Determinants, Cubic Tornheim Jets, and Canonical Digamma
Integrals: Exact harmonic-tail reconstruction, polylogarithmic boundary
identities, and finite-part reduction through the seventh power}.
Repository working report, 11 October 2026 (UTC).
\href{https://raw.githubusercontent.com/VladimirReshetnikov/ProveIt/446dca9bd65c20e409d1d2561d4ffad408103e2a/docs/incoming/ProveIt_Bernoulli_Tornheim_Digamma_2026-10-11.zip}
{\nolinkurl{ProveIt_Bernoulli_Tornheim_Digamma_2026-10-11.zip}}.

\bibitem{CTD}
\emph{Cubic Invariant Tail Descent: All-degree harmonic identities,
exact digamma cubes, polylogarithmic representations, and algebraic
Stieltjes traces}.
Repository working report, 10 October 2026.
\href{https://raw.githubusercontent.com/VladimirReshetnikov/ProveIt/446dca9bd65c20e409d1d2561d4ffad408103e2a/docs/incoming/ProveIt_Cubic_Tail_Descent_2026-10-10.zip}
{\nolinkurl{ProveIt_Cubic_Tail_Descent_2026-10-10.zip}}.

\bibitem{DHR}
\emph{Determinantal Harmonic Resonance: Gauss--Vandermonde finite
parts, Stieltjes jets, and depth-two polylogarithmic closure}.
Repository working report, 10 October 2026.
\href{https://raw.githubusercontent.com/VladimirReshetnikov/ProveIt/446dca9bd65c20e409d1d2561d4ffad408103e2a/docs/incoming/ProveIt_Determinantal_Harmonic_Resonance_2026-10-10.zip}
{\nolinkurl{ProveIt_Determinantal_Harmonic_Resonance_2026-10-10.zip}}.

\bibitem{DPL}
\emph{Differential Polygamma and Lerch Identities: Exact residue
cancellation, higher Gamma moments, spectral derivatives, and
harmonic-tail reconstruction}.
Repository working report, 11 October 2026.
\href{https://raw.githubusercontent.com/VladimirReshetnikov/ProveIt/446dca9bd65c20e409d1d2561d4ffad408103e2a/docs/incoming/ProveIt_Differential_Polygamma_and_Lerch_Jets_2026-10-11.zip}
{\nolinkurl{ProveIt_Differential_Polygamma_and_Lerch_Jets_2026-10-11.zip}}.

\bibitem{MR}
\emph{Contact-Complete Multishift Lerch Reflection: Independent
spectral jets, generalized Stieltjes log moments, confluent contacts,
and Gamma--polylogarithm identities}.
Repository working report, 10 October 2026.
\href{https://raw.githubusercontent.com/VladimirReshetnikov/ProveIt/446dca9bd65c20e409d1d2561d4ffad408103e2a/docs/incoming/ProveIt_Multishift_Reflection_2026-10-10.zip}
{\nolinkurl{ProveIt_Multishift_Reflection_2026-10-10.zip}}.

\bibitem{PTI}
\emph{Exact Polylogarithm Parity, Harmonic-Tail Reconstruction,
and Higher Polygamma Moments: Conjecture proofs, canonical reductions,
spectral derivatives, and a source-level research audit}.
Repository working report, 11 October 2026.
\href{https://raw.githubusercontent.com/VladimirReshetnikov/ProveIt/446dca9bd65c20e409d1d2561d4ffad408103e2a/docs/incoming/ProveIt_Polygamma_Tail_Identities_2026-10-11.zip}
{\nolinkurl{ProveIt_Polygamma_Tail_Identities_2026-10-11.zip}}.

\bibitem{Guo2014}
Li Guo, Georg Regensburger, and Markus Rosenkranz,
\emph{On integro-differential algebras},
Journal of Pure and Applied Algebra \textbf{218}, no.~3 (2014),
456--473.
\url{https://doi.org/10.1016/j.jpaa.2013.06.015}.
Preprint: \url{https://arxiv.org/abs/1212.0266}.

\bibitem{DLMF}
National Institute of Standards and Technology,
\emph{NIST Digital Library of Mathematical Functions}.
Version 1.2.8, released 15 September 2026; consulted 11 October 2026.
In particular, \S5.7 (local Gamma series), \S5.11 (Gamma ratios),
\S5.15 (polygamma functions), \S15.4.20 (Gauss summation),
\S24.8.1 (Bernoulli Fourier series), and \S25.11 (Hurwitz zeta).
\url{https://dlmf.nist.gov/5.7};
\url{https://dlmf.nist.gov/5.11.E13};
\url{https://dlmf.nist.gov/5.15};
\url{https://dlmf.nist.gov/15.4.E20};
\url{https://dlmf.nist.gov/24.8.E1};
\url{https://dlmf.nist.gov/25.11}.

\end{thebibliography}

\end{document}
